%%%%%%%%%%%%%%%%%%%%%%%%%%%%%%%%%%%%%%%%%%%%%%%%%%%%%%%%%%%%%%%%%%%%%%
% Lớp 10 | Mệnh đề. Tập hợp | Mệnh đề
% Dạng 0D1?1-4: Mệnh đề kéo theo, mệnh đề đảo, mệnh đề tương đương
% Số câu: 114  (sinh tự động bởi tach_cau.py; từ nay thầy sửa trực tiếp file này)
%%%%%%%%%%%%%%%%%%%%%%%%%%%%%%%%%%%%%%%%%%%%%%%%%%%%%%%%%%%%%%%%%%%%%%

%%%%%%%%%%%%%%%% TRẮC NGHIỆM 4 PHƯƠNG ÁN %%%%%%%%%%%%%%%%

% nguồn: MaHoaNH_TracNghiem_Lop10
\begin{ex}%[0D1N1-4]%[ma:4aa5ef02]
	Cho mệnh đề \lq\lq$A\Rightarrow B$\rq\rq. Phát biểu nào sau đây không thay thế cho  phát biểu trên?
	\choice  {Nếu $A$ thì $B$}{\True Nếu $B$ thì $A$}{$A$ suy ra $B$}{$A$ là điều kiện đủ để có $B$}
	 \loigiai{
		Mệnh  đề đảo \lq\lq Nếu $B$ thì $A$\rq\rq~ không thay thế cho mệnh đề gốc \lq\lq$A\Rightarrow B$\rq\rq.
	} 
\end{ex}

% nguồn: MaHoaNH_TracNghiem_Lop10
\begin{ex}%[0D1N1-4]%[ma:e55a8c52]
	Trong các  mệnh đề sau, mệnh đề nào là mệnh đề đảo của mệnh đề “Nếu tứ giác $ABCD$ là hình bình hành thì tứ giác $ABCD$ có các góc đối bằng nhau”?
	\choice 
		{Tứ giác $ABCD$ là hình bình hành suy ra tứ giác $ABCD$  có các góc đối bằng nhau}
		{\True Nếu tứ giác $ABCD$ có các góc đối bằng nhau  thì tứ giác $ABCD$ là hình bình hành}
		{Tứ giác $ABCD$ là  hình bình hành nên tứ giác $ABCD$ có các góc đối bằng nhau}
		{Nếu tứ giác $ABCD$ là hình bình hành thì tứ giác $ABCD$ có các góc đối bằng nhau} 
	 \loigiai{
	Mệnh đề đảo của mệnh đề $P\Rightarrow Q$ là mệnh đề $Q\Rightarrow  P$.}
\end{ex}

% nguồn: MaHoaNH_TracNghiem_Lop10
\begin{ex}%[0D1H1-4]%[ma:5dde690d]
	Trong các mệnh đề nào sau đây  mệnh đề nào {\bf sai}?
	\choice 
		{\True Hai tam giác bằng nhau khi và chỉ khi chúng đồng dạng và có  một góc bằng nhau}
		{Một tứ giác là hình chữ nhật khi và chỉ khi chúng  có $3$ góc vuông}
		{Một tam giác là vuông khi và chỉ khi nó có một góc bằng tổng hai góc còn  lại}
		{Một tam giác là đều khi và chỉ khi chúng có hai đường trung tuyến bằng nhau và có một góc  bằng $60^{\circ}$}
	\loigiai{ Hai  tam giác đồng dạng thì chưa chắc bằng nhau.
	} 
\end{ex}

% nguồn: MaHoaNH_TracNghiem_Lop10
\begin{ex}%[0D1H1-4]%[ma:d4733fd8]
	Cho các mệnh đề sau, mệnh  đề nào có mệnh đề đảo đúng?
	 \choice
		{Nếu tứ  giác $ABCD$ là hình thang cân thì $2$ góc đối bù nhau}
		 {Nếu $a=b$ thì $a\cdot c=b\cdot c$}
		 {Nếu $a>b$ thì $a^{2}>b^{2}$}
		{\True Nếu số nguyên chia hết cho $10$ thì  chia hết cho $5$ và $2$}
	 \loigiai{
		\lq\lq Nếu số nguyên chia hết cho $10$ thì chia hết cho $5$ và $2$\rq\rq có mệnh đề đảo là \lq\lq Nếu số nguyên chia hết cho $5$ và $2$ thì chia hết cho $10$\rq\rq là một mệnh đề  đúng.
	 }
\end{ex}

% nguồn: MaHoaNH_TracNghiem_Lop10
\begin{ex}%[0D1H1-4]%[ma:5859a673]
	Cho định lí: \lq\lq Nếu tam giác $ABC$ vuông cân thì tam giác $ABC$ có một góc bằng $45^{\circ}$.\rq\rq Trong các phát biểu sau, phát  biểu nào là mệnh đề đúng?
	 \choice
		{\True Tam giác $ABC$ vuông cân là điều kiện đủ  để tam giác $ABC$ có một góc bằng $45^{\circ}$}
		{Tam giác $ABC$ vuông cân là  điều kiện cần và đủ để tam giác $ABC$ có một góc bằng $45^{\circ}$}
		{Tam giác $ABC$ vuông cân là điều kiện cần để tam giác $ABC$  có một góc bằng $45^{\circ}$}
		{Tam giác $ABC$ có một góc bằng $45^{\circ}$ là điều  kiện đủ để tam giác $ABC$ vuông cân}
	\loigiai{ 
		Tam giác $ABC$ vuông cân là điều kiện đủ để tam giác $ABC$ có một góc  bằng $45^{\circ}$.
	} 
\end{ex}

% nguồn: MaHoaNH_TracNghiem_Lop10
\begin{ex}%[0D1H1-4]%[ma:84d6d30e]
	 Cho các mệnh đề $P$: \lq\lq Hình bình hành $ABCD$ có một góc vuông\rq\rq , $Q$: \lq\lq $ABCD$ là hình chữ nhật\rq\rq. Mệnh đề \lq\lq $P \Rightarrow Q$\rq\rq\, được phát biểu là
	 \choice
		 {Nếu tứ giác $ABCD$ là hình chữ nhật thì $ABCD$ là hình bình hành và có một góc vuông}
		{\True Nếu  hình bình hành $ABCD$ có một góc vuông thì $ABCD$ là hình chữ nhật}
		 {Hình bình hành $ABCD$ có một góc vuông khi và chỉ khi $ABCD$ là hình chữ nhật}
		{Hình bình hành $ABCD$ có một góc vuông là điều kiện cần và đủ để $ABCD$  là hình chữ nhật}
	 \loigiai{
		Mệnh đề \lq\lq  $P \Rightarrow Q$\rq\rq\, được phát biểu là \lq\lq Nếu hình bình hành $ABCD$ có một góc vuông thì $ABCD$ là hình chữ nhật\rq\rq .
	} 
\end{ex}

% nguồn: MaHoaNH_TracNghiem_Lop10
\begin{ex}%[0D1H1-4]%[ma:3b15f856]
	Cho mệnh đề $P \Rightarrow Q$. Khẳng định nào sau đây là  đúng?
	\choice 
		{$Q$ là điều kiện  cần và đủ để có $P$}
		{$P$ là điều kiện cần để  có $Q$}
		{$Q$ là điều kiện đủ  để có $P$}
		{\True $P$ là điều kiện đủ  để có $Q$}
	 \loigiai{
		Ta có  $P$ là điều kiện đủ để có $Q$, $Q$ là điều kiện cần để có $P$.
	} 
\end{ex}

% nguồn: MaHoaNH_TracNghiem_Lop10
\begin{ex}%[0D1H1-4]%[ma:b10320bc]
	Phát biểu mệnh đề đảo của  mệnh đề \lq\lq Nếu $12$ chia hết cho $6$ thì $12$ chia hết cho $3$\rq\rq .
	\choice 
		{Nếu $12$  không chia hết cho $6$ thì $12$ không chia hết cho $3$}
		{\True  Nếu $12$ chia hết cho $3$ thì $12$ chia hết cho $6$}
		{$12$ chia hết cho  $6$ là điều kiện đủ để $12$ chia hết cho $3$}
		{$12$ chia hết cho $6$ khi và chỉ khi 12 chia  hết cho $3$}
	\loigiai{ 
		Mệnh đề đảo của mệnh đề \lq\lq Nếu $12$ chia hết cho $6$ thì $12$ chia hết  cho $3$ là \lq\lq Nếu $12$ chia hết cho $3$ thì $12$ chia hết cho $6$\rq\rq.
	 }
\end{ex}

% nguồn: MaHoaNH_TracNghiem_Lop10
\begin{ex}%[0D1H1-4]%[ma:f1151423]
	Cho hai mệnh đề $P$ và  $Q$. Tìm điều kiện để mệnh đề $P \Rightarrow Q$ là mệnh đề sai.
	 \choice
		{$P$  sai và $Q$ đúng}
		{\True $P$ đúng và  $Q$ sai}
		{$P$ sai và $Q$ sai} 
		{$P$  đúng và $Q$ đúng}
	 \loigiai{
		Ta  có mệnh $P \Rightarrow Q$ chỉ sai khi $P$ đúng $Q$ sai.
	 }
\end{ex}

% nguồn: MaHoaNH_TracNghiem_Lop10
\begin{ex}%[0D1H1-4]%[ma:1472a723]
	 Cho $P \Leftrightarrow Q$ là mệnh đề đúng. Khẳng định nào sau đây là sai?
	\choice 
		{$\overline{P} \Leftrightarrow \overline{Q}$ đúng} 
		{$\overline{Q} \Leftrightarrow P$ sai} 
		{\True $\overline{P} \Leftrightarrow \overline{Q}$ sai} 
		{$\overline{P} \Leftrightarrow Q$ sai} 
	 \loigiai{
		\lq\lq $\overline{P} \Leftrightarrow \overline{Q}$ sai\rq\rq là mệnh đề sai. 
	 }
\end{ex}

% nguồn: MaHoaNH_TracNghiem_Lop10
\begin{ex}%[0D1H1-4]%[ma:c7176f86]
	 Mệnh đề đảo của mệnh đề nào sau đây là một mệnh đề đúng?
	\choice 
		{Nếu một tứ giác là hình thoi thì tứ giác đó có hai đường  chéo vuông góc}
		{\True Nếu một tứ giác có hai đường chéo cắt nhau tại trung điểm mỗi đường thì tứ giác đó là  hình bình hành}
		{Nếu  một hình thang là hình thang cân thì hình thang đó có hai cạnh bên bằng nhau}
		{Nếu một tứ giác là hình vuông thì tứ giác đó có $4$ cạnh bằng nhau} 
	 \loigiai{
		Nếu một tứ giác có hai đường chéo cắt nhau tại trung điểm mỗi đường thì tứ giác đó là hình bình hành có mệnh đề đảo là  mệnh đề đúng.
	 }
\end{ex}

% nguồn: MaHoaNH_TracNghiem_Lop10
\begin{ex}%[0D1H1-4]%[ma:fe9d9e80]
	Mệnh đề nào sau  đây là mệnh đề sai?
	 \choice
		{Hai tam giác có diện tích bằng nhau là  điều kiện cần để chúng là hai tam giác bằng nhau}
		{Nếu hai tam giác bằng nhau thì chúng có diện tích  bằng nhau}
		{\True Hai tam giác bằng nhau là điều kiện cần và đủ để chúng có  diện tích bằng nhau}
		{Hai tam giác bằng nhau  là điều kiện đủ để chúng có diện tích bằng nhau}
	 \loigiai{
		\rq\rq Hai tam giác bằng nhau là điều kiện cần và đủ để chúng có  diện tích bằng nhau là mệnh đề sai.
	 }
\end{ex}

% nguồn: MaHoaNH_TracNghiem_Lop10
\begin{ex}%[0D1H1-4]%[ma:978ff47c]
	 Mệnh đề nào sau đây, có mệnh đề đảo là đúng?
	 \choice
		{Một số tự  nhiên có tận cùng bằng $5$ thì số đó chia hết cho $5$}
		{Nếu  hai tam giác bằng nhau thì diện tích bằng nhau}
		{Nếu  $a$ và $b$ cùng chia hết cho $c$ thì $ab$ chia hết cho $c$}
		{\True Nếu $a$ chia hết  cho $2$ thì $a+1$ là số lẻ}
	 \loigiai{
		Nếu $a+1$ là  số lẻ thì $a$ là số chẵn nên $a$ sẽ chia hết cho $2$. Do đó mệnh đề \lq\lq Nếu $a$ chia hết cho $2$ thì $a+1$ là số lẻ\rq\rq là mệnh đề có mệnh đề đảo đúng.
	} 
\end{ex}

% nguồn: MaHoaNH_TracNghiem_Lop10
\begin{ex}%[0D1H1-4]%[ma:fade67c3]
	Trong các mệnh đề sau, mệnh đề  nào là mệnh đề \textbf{sai}?
	\choice 
		{$\sqrt{23}<5 \Rightarrow  -2\sqrt{23}>-2 \cdot 5$}
		{$\pi<4 \Leftrightarrow  \pi^{2}<16$}
		 {\True $-\pi<-2 \Leftrightarrow \pi^{2}<4$}
		 {$\sqrt{23}<5 \Rightarrow 2\sqrt{23}<2.5$}
	 \loigiai{
		Do $\pi^{2}<4$ là sai nên  mệnh đề $-\pi<-2 \Leftrightarrow \pi^{2}<4$ là sai.
	} 
\end{ex}

% nguồn: MaHoaNH_TracNghiem_Lop10
\begin{ex}%[0D1H1-4]%[ma:130c966e]
	Cho mệnh đề $P \Rightarrow Q$  sai khi
	\choice 
		{$P$ đúng và  $Q$ đúng}
		{$P$ sai  và $Q$ đúng}
		{$P$ sai và $Q$ sai} 
		{\True $P$ đúng và $Q$ sai} 
	 \loigiai{
		Mệnh  đề $P \Rightarrow Q$ sai khi $P$ đúng và $Q$ sai.
	} 
\end{ex}

% nguồn: MaHoaNH_TracNghiem_Lop10
\begin{ex}%[0D1H1-4]%[ma:65faa958]
	Mệnh  đề nào sau đây {\bf sai}?
	\choice 
		{Tứ giác  $ABCD$ là hình chữ nhật $\Rightarrow$ tứ giác $ABCD$ có ba góc vuông}
		{\True Tam giác $ABC$ là  tam giác đều $\Leftrightarrow$ $\widehat{A}=60^\circ$}
		{Tam  giác $ABC$ cân tại $A$ $\Rightarrow$ $AB=AC$}
		{Tứ giác  $ABCD$ nội tiếp đường tròn tâm $O$ $\Rightarrow$ $OA=OB=OC=OD$}
	\loigiai{ 
		Tam giác $ABC$ là tam giác đều $\Leftrightarrow$  $\widehat{A}=60^\circ$ là hai mệnh đề không tương đương, tam giác $ABC$ có $\widehat{A}=60^\circ$ chưa đủ để nó là tam giác đều.
	 }
\end{ex}

% nguồn: MaHoaNH_TracNghiem_Lop10
\begin{ex}%[0D1H1-4]%[ma:749c1107]
	Cho tam giác $ABC$ với $H$ là chân đường cao từ $A$. Mệnh đề nào sau  đây {\bf sai}?
	\choice 
		{\lq\lq $ABC$ là tam giác vuông ở $A$ $\Leftrightarrow \dfrac{1}{AH^2}=\dfrac{1}{AB^2}+\dfrac{1}{AC^2}$\rq\rq} 
		{\lq\lq  $ABC$ là tam giác vuông ở $A$ $\Leftrightarrow BA^2=BH\cdot BC$\rq\rq}
		{\lq\lq  $ABC$ là tam giác vuông ở $A$ $\Leftrightarrow HA^2=HB\cdot HC$\rq\rq}
		{\True \lq\lq $ABC$ là tam giác vuông ở $A$ $\Leftrightarrow BA^2=BC^2+AC^2$\rq\rq} 
	 \loigiai{
		Chọn đúng phải là \lq\lq $ABC$ là tam  giác vuông ở $A$ $\Leftrightarrow BC^2=AB^2+AC^2$\rq\rq.
	 }
\end{ex}

% nguồn: MaHoaNH_TracNghiem_Lop10
\begin{ex}%[0D1H1-4]%[ma:ec0ab31a]
	Trong các mệnh đề sau, mệnh đề  nào có mệnh đề đảo đúng?
	\choice 
		{Nếu số nguyên $n$ có chữ số tận  cùng là $5$ thì số nguyên $n$ chia hết cho $5$}
		{\True Nếu tứ giác $ABCD$ có hai đường chéo cắt  nhau tại trung điểm mỗi đường thì tứ giác $ABCD$ là hình bình hành}
		{Nếu tứ giác $ABCD$ là hình chữ nhật thì tứ giác $ABCD$ có  hai đường chéo bằng nhau}
		{Nếu tứ giác $ABCD$ là hình thoi thì tứ giác $ABCD$ có hai đường chéo  vuông góc với nhau}
	\loigiai{ 
		Xét mệnh đề đảo của mệnh đề: \lq\lq Nếu tứ giác $ABCD$ là hình bình hành thì tứ giác $ABCD$ có hai đường chéo cắt nhau tại trung điểm mỗi đường\rq\rq . Mệnh  đề này đúng.
	 }
\end{ex}

% nguồn: MaHoaNH_TracNghiem_Lop10
\begin{ex}%[0D1H1-4]%[ma:db1552f2]
	Tìm mệnh đề đảo  của mệnh đề \lq\lq Nếu tam giác có $2$ cạnh bằng nhau thì tam giác đó là tam giác cân\rq\rq .
	\choice 
		{Tam giác  đó là tam giác cân}
		{Một tam giác là tam giác cân nếu và chỉ nếu tam giác đó có $2$ cạnh bằng nhau} 
		{Một tam giác không có hai cạnh bằng nhau thì tam giác  đó không là tam giác cân}
		{\True Nếu một tam giác là tam giác cân thì tam giác đó  có hai cạnh bằng nhau}
	\loigiai{ 
		Mệnh đề đảo của mệnh  đề, \lq\lq Nếu tam giác có $2$ cạnh bằng nhau thì tam giác đó là tam giác cân\rq\rq là \lq\lq Nếu một tam giác là tam giác cân thì tam giác đó có hai cạnh bằng nhau\rq\rq .
	} 
\end{ex}

% nguồn: MaHoaNH_TracNghiem_Lop10
\begin{ex}%[0D1H1-4]%[ma:d64a8e12]
	Cho định lý \lq\lq Nếu hai  tam giác bằng nhau thì diện tích bằng nhau\rq\rq. Mệnh đề nào sau đây \textbf{đúng}?
	\choice 
		{Hai tam giác bằng  nhau là điều kiện cần và đủ để chúng có diện tích bằng nhau}
		{Hai tam giác bằng  nhau là điều kiện cần để diện tích chúng bằng nhau}
		{Hai tam giác có diện tích bằng nhau là điều kiện  đủ để chúng bằng nhau}
		{\True Hai tam giác  bằng nhau là điều kiện đủ để diện tích chúng bằng nhau}
	\loigiai{ 
		Hai tam giác bằng nhau là điều kiện đủ để  diện tích chúng bằng nhau.
	 }
\end{ex}

% nguồn: MaHoaNH_TracNghiem_Lop10
\begin{ex}%[0D1H1-4]%[ma:91b8b6a6]
	Cho mệnh đề \rq\rq $\forall x \in X$ $P(x) \Rightarrow Q(x)$\rq\rq . Chọn khẳng định không đúng. 
	 \choice
		{$P(x)$ là điều kiện đủ để có  $Q(x)$}
		{$Q(x)$ là điều kiện  cần để có $P(x)$}
		{$P(x)$ là giả thiết  và $Q(x)$ là kết luận}
		{\True $P(x)$ là  điều kiện cần để có $Q(x)$}
	\loigiai{ 
		Ta  có $P(x)$ là điều kiện cần để có $Q(x)$.
	 }
\end{ex}

% nguồn: MaHoaNH_TracNghiem_Lop10
\begin{ex}%[0D1H1-4]%[ma:cfac5f9f]
	Trong các mệnh đề sau, mệnh đề nào là mệnh đề sai? 
	 \choice
		{\lq\lq $ABC$ là tam giác đều $\Leftrightarrow$ Tam giác  $ABC$ cân\rq\rq }
		 {\lq\lq $ABC$ là tam giác đều $\Leftrightarrow$ Tam giác $ABC$ cân và có một góc $60^{\circ}$\rq\rq }
		{\lq\lq $ABC$ là tam giác đều $\Leftrightarrow$ $ABC$ là tam giác có  ba cạnh bằng nhau\rq\rq }
		{\True  \lq\lq $ABC$ là tam giác đều $\Leftrightarrow$ Tam giác $ABC$ có hai góc bằng $60^{\circ}$\rq\rq }
	\loigiai{ 
		Mệnh đề kéo theo \lq\lq $ABC$ là tam giác đều $\Leftrightarrow$ Tam giác $ABC$ cân\rq\rq là mệnh đề đúng, nhưng mệnh đề đảo \lq\lq  Tam giác $ABC$ cân $\Leftrightarrow$ $ABC$ là tam giác đều\rq\rq là mệnh đề sai. Do đó, $2$ mệnh đề \lq\lq $ABC$ là tam giác đều\rq\rq và \lq\lq Tam giác $ABC$ cân\rq\rq không phải là $2$ mệnh đề tương đương.
	} 
\end{ex}

% nguồn: MaHoaNH_TracNghiem_Lop10
\begin{ex}%[0D1H1-4]%[ma:35ed614a]
	Cho mệnh đề $P$: \lq\lq Nếu hai tam giác bằng nhau thì diện tích của chúng bằng nhau\rq\rq. Một  phát biểu khác của mệnh đề $P$ là
	 \choice
		{Hai tam giác bằng  nhau là điều kiện cần để diện tích chúng bằng nhau}
		{Hai tam giác bằng nhau là điều kiện cần và đủ để chúng  có diện tích bằng nhau}
		{\True Hai tam giác có diện tích bằng nhau là điều  kiện cần để chúng bằng nhau}
		{Hai tam giác có diện tích bằng nhau là điều  kiện đủ để chúng bằng nhau}
	\loigiai{ 
		Cách phát biểu khác của $P$: \lq\lq Hai tam giác có diện tích bằng nhau là điều  kiện cần để chúng bằng nhau\rq\rq.
	 }
\end{ex}

% nguồn: MaHoaNH_TracNghiem_Lop10
\begin{ex}%[0D1H1-4]%[ma:08bd4aca]
	Trong các mệnh đề sau, mệnh đề nào {\bf  sai}?
	 \choice
		{\True $-\pi <-2 \Leftrightarrow  \pi^2 < 4$}
		{$\pi < 4 \Leftrightarrow \pi^2 <  16$}
		{$\sqrt{23} < 5 \Rightarrow 2\cdot \sqrt{23}  < 2\cdot 5$}
		{$\sqrt{23} < 5  \Rightarrow-2\cdot \sqrt{23} >-2\cdot 5$}
	 \loigiai{
		 Mệnh đề kéo theo chỉ sai khi $P$ đúng, $Q$ sai.\\
		Vậy mệnh đề  ở đáp án $-\pi <-2 \Leftrightarrow \pi^2 < 4$ sai.
	 }
\end{ex}

% nguồn: MaHoaNH_TracNghiem_Lop10
\begin{ex}%[0D1H1-4]%[ma:3c6dac9c]
	Trong các mệnh đề sau, mệnh  đề nào có mệnh đề đảo đúng?
	\choice 
		 {Nếu số nguyên $n$ có tổng các chữ số bằng $9$ thì số tự nhiên $n$ chia hết cho $3$}
		{Nếu $x>y$ thì  $x^{2}>y^{2}$}
		 {Nếu $x=y$ thì $t\cdot x=t\cdot y$}
		{\True Nếu $x>y$ thì $x^{3}>y^{3}$} 
	\loigiai{ 
		Xét mệnh đề  đảo của mệnh đề \lq\lq Nếu số tự nhiên $n$ chia hết cho $3$ thì số nguyên $n$ có tổng các chữ số bằng $9$\rq\rq . Mệnh đề này sai vì tổng các chữ số của $n$ phải chia hết cho $9$ thì $n$ mới chia hết cho $9$.\\
		Xét mệnh đề đảo của mệnh đề \lq\lq Nếu $x^{2}>y^{2}$ thì $x>y$\rq\rq sai  vì $x^{2}>y^{2}\Leftrightarrow|x|>|y|\Leftrightarrow\hoac{&x>y\\ &x<-y.}$\\
		Xét mệnh đề đảo của mệnh đề \lq\lq Nếu $t\cdot x=t\cdot y$  thì $x=y$\rq\rq sai với $t=0\Rightarrow x,$ $y\in\mathbb{R}$.
	} 
\end{ex}

% nguồn: MaHoaNH_TracNghiem_Lop10
\begin{ex}%[0D1H1-4]%[ma:bbef1c67]
	Cho mệnh đề $P$: \lq\lq Hai số nguyên chia hết cho $7$\rq\rq\, và mệnh đề $Q$: \lq\lq Tổng của chúng chia hết cho $7$\rq\rq . Phát  biểu mệnh đề $P \Rightarrow Q$.
	\choice 
		{Nếu hai số nguyên chia hết cho $7$ thì tổng của chúng không  chia hết cho $7$}
		{\True Nếu hai số  nguyên chia hết cho $7$ thì tổng của chúng chia hết cho $7$}
		{Nếu hai số nguyên không chia hết cho  $7$ thì tổng của chúng không chia hết cho 7}
		{Nếu tổng của hai số nguyên chia hết cho $7$ thì hai số nguyên  đó chia hết cho 7}
	 \loigiai{
		Mệnh đề $P$: \lq\lq Hai  số nguyên chia hết cho $7$\rq\rq.\\
		Mệnh đề $Q$: \lq\lq Tổng của chúng  chia hết cho $7$\rq\rq.\\
		Mệnh đề $P \Rightarrow  Q$ có dạng: \lq\lq Nếu $P$ thì $Q$\rq\rq.\\
		Vậy mệnh đề $P \Rightarrow Q$: \lq\lq Nếu hai  số nguyên chia hết cho $7$ thì tổng của chúng chia hết cho $7$\rq\rq .
	} 
\end{ex}

% nguồn: MaHoaNH_TracNghiem_Lop10
\begin{ex}%[0D1B1-4]%[ma:dc2da9b7]
	Trong các câu  sau, câu nào là mệnh đề đúng?
	 \choice
		{Nếu $a\ge b$ thì $a^2\geq b^2$} 
		{\True Nếu $a$ chia hết cho  $9$ thì $a$ chia hết cho $3$}
		{Nếu $a$ chia hết cho $5$ thì  $a$ có chữ số tận cùng là $5$}
		 {Nếu một tam giác có một góc bằng $60^\circ $ thì tam giác đó đều}
	\loigiai{Vì $9$ chia hết cho  $3$ nên nếu $a$ chia hết cho $9$ thì $a$ chia hết cho $3$. }
\end{ex}

% nguồn: MaHoaNH_TracNghiem_Lop10
\begin{ex}%[0D1B1-4]%[ma:3dc68b04]
	Mệnh đề nào sau đây là mệnh đề \textbf{sai}? 
	 \choice
		{Một tam giác là đều khi và chỉ khi chúng có hai đường trung tuyến bằng nhau và  có một góc bằng $60^\circ$}
		{Một tam giác là vuông khi và  chỉ khi nó có một cạnh bình phương bằng tổng bình phương hai cạnh còn lại}
		{Một tứ giác là hình  chữ nhật khi và chỉ khi chúng có $3$ góc vuông}
		{\True Hai tam giác bằng nhau khi và chỉ khi  chúng đồng dạng và có một góc bằng nhau}
	\loigiai{Phương án: \lq\lq Hai tam giác bằng nhau khi và chỉ khi chúng đồng dạng và  có một góc bằng nhau\rq\rq \ là sai. Ví dụ: Hai tam giác đều có độ dài cạnh bằng $1$ và $2$ đồng dạng nhưng không bằng nhau.}
\end{ex}

% nguồn: MaHoaNH_TracNghiem_Lop10
\begin{ex}%[0D1B1-4]%[ma:6fb57d82]
	Trong các mệnh đề  sau đây, mệnh đề nào đúng?
	\choice 
		{\True Với mọi số thực $x$, nếu  $x<-2$ thì $x^2>4$}
		{Với  mọi số thực $x$, nếu $x^2<4$ thì $x<-2$}
		{Với mọi số thực  $x$, nếu $x<-2$ thì $x^2<4$}
		{Với mọi số thực $x$, nếu $x^2>4$  thì $x>-2$}
	 \loigiai{
	Ta thấy  $x<-2<0\Rightarrow -x>2>0\Rightarrow x^2>4$ là đúng. }
\end{ex}

% nguồn: MaHoaNH_TracNghiem_Lop10
\begin{ex}%[0D1B1-4]%[ma:b7a72a61]
	Trong các mệnh đề sau, mệnh đề nào  là mệnh đề \textbf{sai}?
	 \choice
		{\True \lq\lq $\triangle ABC$ là tam giác đều $\Leftrightarrow  $ tam giác $ABC$ cân\rq\rq}
		{\lq\lq $\triangle ABC$ là tam giác đều $\Leftrightarrow $ Tam giác $ABC$ cân  và có một góc $60^\circ$\rq\rq}
		{\lq\lq $\triangle ABC$ là tam giác đều $\Leftrightarrow \triangle ABC$ là tam giác có ba cạnh bằng nhau \rq\rq} 
		{\lq\lq $ABC$ là tam giác đều $\Leftrightarrow$  tam giác $\triangle ABC$ có hai góc bằng nhau \rq\rq}
	\loigiai{ 
	Tam giác $ABC$ cân cần thêm các yếu tố khác  thì mới có thể trở thành tam giác đều.}
\end{ex}

% nguồn: MaHoaNH_TracNghiem_Lop10
\begin{ex}%[0D1B1-4]%[ma:da30326d]
	Trong các mệnh đề sau, mệnh đề nào  là mệnh đề \textbf{sai}?
	\choice 
		{$-\pi<-2\Leftrightarrow{{\pi}^2}<4$} 
		{$\pi<4\Leftrightarrow{{\pi}^2}<16$} 
		 {$\sqrt{23}<5\Rightarrow 2\sqrt{23}<2\text{,}5$}
		{\True  $\sqrt{23}<5\Rightarrow-2\sqrt{23}>-2\text{,}5$}
	 \loigiai{
	Ta thấy $\sqrt{23}<5$ là đúng nhưng $-2\sqrt{23}>-2\text{,}5$ là sai nên  \lq\lq $\sqrt{23}<5\Rightarrow-2\sqrt{23}>-2\text{,}5$\rq\rq\ là sai. }
\end{ex}

% nguồn: MaHoaNH_TracNghiem_Lop10
\begin{ex}%[0D1K1-4]%[ma:72f5baef]
	Cho $A, B$ và $C$ là các mệnh đề. Biết  rằng các mệnh đề $A, B$ và $A \Rightarrow (B \Rightarrow \overline{C})$ là các mệnh đề đúng. Phát biểu nào sau đây là đúng?
	\choice 
		{$A \Rightarrow \overline{B}$ là mệnh  đề đúng}
		{\True $A \Rightarrow {C}$ là mệnh  đề sai}
		{$A \Leftrightarrow B$  là mệnh đề sai}
		{$\overline{C} \Rightarrow \overline{B}$ là mệnh đề đúng} 
	\loigiai{Do $B$ đúng và $B \Rightarrow \overline{C}$ đúng nên $\overline{C}$  là mệnh đề đúng.\\
	Suy ra $C$ là mệnh đề  sai. Mà $A$ đúng nên mệnh đề $A \Rightarrow {C}$ là mệnh đề sai.}
\end{ex}

% nguồn: MaHoaNH_TracNghiem_Lop10
\begin{ex}%[0D1B1-4]%[ma:c0326c6d]
	Cho định lí “Nếu  hai tam giác bằng nhau thì diện tích chúng bằng nhau”. Mệnh đề nào sau đây đúng?
	\choice 
		{Hai tam giác bằng nhau  là điều kiện cần để diện tích chúng bằng nhau}
		{Hai tam giác bằng  nhau là điều kiện cần và đủ để chúng có diện tích bằng nhau}
		{Hai tam giác có diện tích bằng nhau là  điều kiện đủ để chúng bằng nhau}
		{\True Hai tam giác bằng nhau là điều kiện  đủ để diện tích chúng bằng nhau}
	 \loigiai{
	“Hai tam giác bằng nhau” là điều kiện đủ.“Diện tích bằng nhau” là điều  kiện cần.}
\end{ex}

% nguồn: MaHoaNH_TracNghiem_Lop10
\begin{ex}%[0D1B1-4]%[ma:2ef23c9b]
	Trong các mệnh đề sau, mệnh đề nào có mệnh  đề đảo đúng?
	\choice 
		{Nếu  số nguyên $n$ có chữ số tận cùng là $5$ thì số nguyên $n$ chia hết cho $5$}
		{Nếu tứ giác $ABCD$ có hai đường chéo cắt nhau tại trung điểm mỗi  đường thì tứ giác $ABCD$ là hình bình hành}
		{\True Nếu tứ giác $ABCD$ là hình chữ nhật  thì tứ giác $ABCD$ có hai đường chéo bằng nhau}
		{Nếu tứ giác $ABCD$ là hình thoi thì tứ giác $ABCD$ có hai đường chéo vuông góc  với nhau}
	\loigiai{ 
	Mệnh đề đảo của mệnh đề: \lq\lq Nếu tứ giác $ABCD$ có hai đường chéo cắt nhau tại trung điểm mỗi đường thì tứ giác $ABCD$ là hình bình hành \rq\rq\, là mệnh đề \lq\lq Nếu tứ giác $ABCD$ là hình bình hành  thì nó có hai đường chéo cắt nhau tại trung điểm mỗi đường \rq\rq\, là một mệnh đề đúng.}
\end{ex}

% nguồn: MaHoaNH_TracNghiem_Lop10
\begin{ex}%[0D1B1-4]%[ma:23d2b944]
	Trong các mệnh đề sau, mệnh đề nào  là mệnh đề \textbf{sai}?
	 \choice
		 {\True Hai tam giác bằng nhau khi và chỉ khi chúng đồng dạng và có một góc bằng nhau}
		{Một tứ giác là hình chữ nhật khi và chỉ khi chúng có $3$ góc  vuông}
		{Một tam giác là vuông khi và chỉ  khi nó có một góc bằng tổng hai góc còn lại}
		{Một tam giác là đều khi và chỉ  khi chúng có hai đường trung tuyến bằng nhau và có một góc bằng $60^\circ$}
	\loigiai{ 
	Nếu hai tam giác đồng dạng thì chúng luôn có ba góc tương ứng bằng nhau nhưng chưa đủ cơ sở để khẳng định  chúng bằng nhau. }
\end{ex}

% nguồn: MaHoaNH_TracNghiem_Lop10
\begin{ex}%[0D1B1-4]%[ma:79224716]
	Cho phương trình bậc hai $ax^2+bx+c=0$ có $\Delta=b^2-4ac$. Khẳng định nào sau đây đúng? 
	 \choice
		{\True Nếu  $\Delta\geq 0$ thì phương trình có hai nghiệm phân biệt}
		{Nếu $b^2>4ac$  thì phương trình có hai nghiệm trái dấu}
		{Nếu  $a\cdot c>0$ thì phương trình có hai nghiệm dương}
		{Nếu $a\cdot c<0$ thì phương trình  có hai nghiệm phân biệt}
	\loigiai{ 
	Nếu $a\cdot c<0$ thì $\Delta=b^2-4ac>0$, với mọi $a,\,b,\, c$ nên phương  trình luôn có hai nghiệm phân biệt và là hai nghiệm trái dấu.}
\end{ex}

% nguồn: MaHoaNH_TracNghiem_Lop10
\begin{ex}%[0D1K1-4]%[ma:89ec7de5]
	 Cho hai mệnh đề $P$ và $Q$. Mệnh đề $P\Leftrightarrow Q$ đúng khi
	\choice 
		{$P$ đúng và $Q$  sai}
		{$\overline{P}$ đúng và  $Q$ đúng}
		{$P$ sai  và $Q$ đúng}
		 {\True $\overline{P}$ sai và $\overline{Q}$ sai}
	 \loigiai{
		Ta có $\overline{P}$ sai và $\overline{Q}$ sai, suy ra P  đúng và Q đúng nên $P\Leftrightarrow Q$ đúng.
	 }
\end{ex}

% nguồn: MaHoaNH_TracNghiem_Lop10
\begin{ex}%[0D1K1-4]%[ma:2495b9a8]
	 Cho định lý: \lq\lq Với mọi số tự nhiên $n$ nếu $n^2\ \vdots\ 3$ thì $n\ \vdots\ 3$ \rq\rq. Khẳng định nào sau đây là đúng?
	\choice 
		{Giả thiết  là \lq\lq $n\ \vdots\ 3$\rq\rq}
		{Giả thiết là \lq\lq  nếu $n^2\ \vdots\ 3$\rq\rq}
		{\True Kết luận là \lq\lq $n\ \vdots\ 3$\rq\rq} 
		{Kết luận là \lq\lq thì  $n\ \vdots\ 3$\rq\rq}
	\loigiai{ 
		Giả thiết là \lq\lq $n^2\ \vdots\ 3$\rq\rq, kết luận là \lq\lq $n\ \vdots\ 3$\rq\rq. 
	} 
\end{ex}

% nguồn: MaHoaNH_TracNghiem_Lop10
\begin{ex}%[0D1B1-4]%[ma:fb3095b7]
	Mệnh đề nào sau đây là \textbf{sai}? 
	\choice 
		{Điều  kiện cần và đủ để phương trình bậc hai $ax^2+bx+c=0$ ($ a\ne 0 $) vô nghiệm là $\Delta = b^2-4ac<0$}
		{Số nguyên $n$ chia hết cho $5$ khi và chỉ khi số tận cùng của $n$ phải là số $0$ hoặc số  $5$}
		{\True Điều kiện cần  và đủ để để $\triangle ABC$ đều là $\triangle ABC$ cân}
		{Số nguyên $n$  là số chẵn khi và chỉ khi $n$ chia hết cho $2$}
	\loigiai{ 
		 Mệnh đề ``Điều kiện cần và đủ để để $\triangle ABC$ đều là $\triangle ABC$ cân'' sai vì $\triangle ABC$ cân không đủ suy ra $\triangle ABC$ đều.
	 }
\end{ex}

% nguồn: MaHoaNH_TracNghiem_Lop10
\begin{ex}%[0D1K1-4]%[ma:d7810ecd]
	Trong các mệnh dề sau.  Mệnh đề nào không phải là định lí?
	\choice 
		{Điều kiện đủ để trong mặt phẳng, hai đường thẳng song song với nhau là hai đường thẳng ấy  cùng vuông góc với đường thẳng thứ ba}
		{Điều kiện đủ để diện tích hai tam giác bằng nhau là hai  tam giác ấy bằng nhau}
		{Điều kiện đủ để  hai đường chéo của một tứ giác vuông góc với nhau là tứ giác ấy là hình thoi}
		{\True Điều kiện đủ để một số nguyên dương $a$ tận cùng bằng $5$ là số  đó chia hết cho $5$}
	\loigiai{Đầu tiên ta phân biệt mệnh đề \lq\lq Nếu $P$ thì $Q$\rq\rq thì điều kiện đủ để có $Q$ là $P$ nên ta thấy mệnh đề \lq\lq Điều kiện đủ để một số nguyên dương $a$ tận cùng bằng $5$ là số đó  chia hết cho $5$\rq\rq\, là sai vì số chia hết cho $5$ chưa hẵn có tận cùng là $5$.\\
		Vậy khẳng định  như vậy không phải là một định lí.
	 }
\end{ex}

% nguồn: MaHoaNH_TracNghiem_Lop10
\begin{ex}%[0D1B1-4]%[ma:3df91739]
	Với  số thực $x$ bất kì, mệnh đề nào sau đây đúng?
	 \choice
		 {\True $\forall x,x^2>5\Rightarrow x>\sqrt{5} $ hoặc $x<-\sqrt{5}$}
		{$  \forall x,x^2>5\Rightarrow -\sqrt{5}<x<\sqrt{5}$}
		{$\forall  x,x^2>5\Rightarrow x>\pm \sqrt{5}$}
		{$\forall x,x^2>5\Rightarrow x\ge  \sqrt{5}$ hoặc $x\leq -\sqrt{5}$}
	\loigiai{ 
		 Ta có $\forall x,x^2>5\Rightarrow x>\sqrt{5} $ hoặc $x<-\sqrt{5}$.
	} 
\end{ex}

% nguồn: MaHoaNH_TracNghiem_Lop10
\begin{ex}%[0D1B1-4]%[ma:843996ed]
	Xét trên tập hợp số tự nhiên, mệnh đề nào là  \textbf{sai}?
	 \choice
		{Nếu một số chia hết cho $2$ thì có chữ số tận cùng  là chữ số chẵn}
		{\True Nếu thì một số chia hết cho $5$  thì nó có chữ số tận cùng là $5$}
		{Số có tổng các chữ số chia hết cho $3$ thì chia hết cho  $3$}
		 {Số chẵn chia hết cho $3$ thì chia hết cho $6$}
	 \loigiai{
		Nếu thì một số chia hết cho $5$ thì nó có chữ số tận cùng là $5$ hoặc $0$. 
	 }
\end{ex}

% nguồn: MaHoaNH_TracNghiem_Lop10
\begin{ex}%[0D1K1-4]%[ma:2d151cfe]
	Biết $A$ là mệnh đề đúng, $B$ là mệnh đề sai, $C$ là mệnh đề đúng. Mệnh đề nào  sau đây \textbf{sai}?
	 \choice
		{$A\Rightarrow C$} 
		{$C\Rightarrow (A\Rightarrow  \overline{B})$}
		{$(\overline{B}\Rightarrow C)\Rightarrow A$} 
		{\True  $C\Rightarrow (A\Rightarrow B)$}
	\loigiai{Cần lưu ý rằng  mệnh đề ``$P\Rightarrow Q$'' chỉ sai khi $P$ đúng và $Q$ sai.\\
		 Do đó mệnh đề ``$A\Rightarrow B$'' là mệnh đề sai.\\
	Suy ra mệnh đề ``$C\Rightarrow (A\Rightarrow B)$''  sai.}
\end{ex}

% nguồn: MaHoaNH_TracNghiem_Lop10
\begin{ex}%[0D1B1-4]%[ma:8d32310e]
	Trong các mệnh đề tương đương sau đây, mệnh đề nào \textbf{sai}? 
	\choice 
		{$n $ là số nguyên lẻ khi và khi $n^2$  là số lẻ}
		{$ n$ chia hết cho  $3$ khi và chỉ khi tổng các chữ số của $n$ chia hết cho $3$}
		{\True $ABCD $ là  hình chữ nhật khi và chỉ khi $AC=BD$}
		{$ABC $ là  tam giác đều khi và chỉ khi $AB=AC$ và $\widehat{A}=60^\circ$}
	 \loigiai{
		 Mệnh đề \lq\lq$ABCD $ là hình chữ nhật khi và chỉ khi $AC=BD$\rq\rq\ sai vì khi $AC=BD$ thì $ABCD$ chưa phải là hình chữ nhật.
	 }
\end{ex}

% nguồn: MaHoaNH_TracNghiem_Lop10
\begin{ex}%[0D1B1-4]%[ma:98b65971]
	 Cách phát biểu nào sau đây là sai về mệnh đề $P\Rightarrow Q$
	\choice 
		{$P$ là điều kiện  đủ để có $Q$}
		{$P$ kéo theo $Q$} 
		{\True $Q$ là điều  kiện đủ để có $P$}
		{$Q$ là điều kiện cần  để có $P$}
	\loigiai{ 
		Phát biểu sai: $Q$ là điều kiện đủ để  có $P$.
	 }
\end{ex}

% nguồn: MaHoaNH_TracNghiem_Lop10
\begin{ex}%[0D1B1-4]%[ma:445e313d]
	Trong các mệnh đề sau, mệnh đề nào có mệnh đề  đảo là mệnh đề đúng?
	\choice 
		{Nếu $2$ tam giác bằng nhau thì hai tam giác đó có diện tích bằng  nhau}
		{Nếu một số tự nhiên có chữ số tận cùng bằng $ 0 $ thì số đó chia hết  cho $ 5 $}
		{Nếu $a$ và $b$ chia hết cho $c$ thì $a + b$  chia hết cho $c$}
		{\True Nếu $a$ chia hết cho  $ 3 $ thì $a$ chia hết cho $ 9 $}
	 \loigiai{
		Ta có ``Nếu $a$ chia hết cho $9$ thì $a$ chia hết cho $3$'' là  mệnh đề đúng.
	 }
\end{ex}

% nguồn: MaHoaNH_TracNghiem_Lop10
\begin{ex}%[0D1B1-4]%[ma:39957130]
	Với số thực  $x$ bất kì, mệnh đề nào sau đây đúng?
	\choice 
		{$\forall x,x^2\leq 16\Leftrightarrow  x\leq \pm 4 $}
		{\True $\forall x,x^2\leq 16\Leftrightarrow -4\leq x\leq 4  $}
		{$\forall x,x^2\leq  16\Leftrightarrow x\leq -4, x\ge 4$}
		{$ \forall  x,x^2\leq 16\Leftrightarrow -4<x<4$}
	 \loigiai{
		Ta có $\forall x,x^2\leq 16\Leftrightarrow |x|\le 16  \Leftrightarrow -4\leq x\leq 4 $.
	} 
\end{ex}

% nguồn: MaHoaNH_TracNghiem_Lop10
\begin{ex}%[0D1B1-4]%[ma:358da832]
	Trong các mệnh đề sau, mệnh  đề nào có mệnh đề đảo đúng?
	 \choice
		{Nếu $a $ và $b$ chia hết cho $c$ thì $a+b$ chia hết cho $c$} 
		{Nếu hai  tam giác bằng nhau thì có diện tích bằng nhau}
		{\True Nếu $a $ chia  hết cho $3$ thì $a$ chia hết cho $9$}
		{Nếu một số tận  cùng bằng $0$ thì số đó chia hết cho $5$}
	 \loigiai{
		Mệnh đề: \lq\lq Nếu $a $ chia hết cho $3$ thì $a$ chia hết cho $9$\rq\rq\ có mệnh đề đảo là \lq\lq Nếu $a $ chia hết cho $9$ thì $a$ chia hết cho $3$\rq\rq\ đúng. 
	} 
\end{ex}

% nguồn: MaHoaNH_TracNghiem_Lop10
\begin{ex}%[0D1B1-4]%[ma:dacdd809]
	Mệnh đề nào sau đây là  \textbf{sai}?
	 \choice
		{Tứ  giác $ABCD$ là hình chữ nhật $\Leftrightarrow $ tứ giác $ABCD$ có 3 góc vuông}
		{\True Tam giác $ABC$ đều $\Leftrightarrow \widehat{A}=60^\circ$} 
		{Tam giác $ABC$ cân tại $A\Leftrightarrow AB=AC$} 
		{Một tam giác là tam giác vuông $\Leftrightarrow$ nó có một góc bằng  tổng hai góc còn lại}
	\loigiai{ 
		Mệnh đề ``Tam giác $ABC$ đều $\Leftrightarrow \widehat{A}=60^\circ$'' là mệnh đề sai vì $ \widehat{A}=60^\circ $ không đủ suy ra  tam giác $ABC$ đều.
	 }
\end{ex}

% nguồn: MaHoaNH_TracNghiem_Lop10
\begin{ex}%[0D1B1-4]%[ma:25dbab11]
	Phát biểu nào sau đây sai? 
	\choice 
		{Điều kiện cần và đủ để phương trình bậc  hai $a x^{2}+b x+c=0$ vô nghiệm là $\Delta=b^{2}-4 a c<0$}
		{Số nguyên $n$ chia hết cho $5$ khi và chỉ khi số  tận cùng của $n$ phải là $0$ hoặc $5$}
		{Một tam giác là tam  giác vuông khi và chỉ khi nó có một góc bằng tổng hai góc còn lại}
		{\True  Điều kiện cần và đủ để $\triangle ABC$ đều là $\triangle ABC$ cân}
	 \loigiai{
		$\triangle ABC$ đều suy ra $\triangle ABC$ cân nhưng điều ngược lại nói chung không đúng. Vậy khẳng định  sai là: Điều kiện cần và đủ để $\triangle ABC$ đều là $\triangle ABC$ cân.
	} 
\end{ex}

% nguồn: MaHoaNH_TracNghiem_Lop10
\begin{ex}%[0D1K1-4]%[ma:be5f704e]
	Biết $A$ là  mệnh đề sai, còn $B$ là mệnh đề đúng. Mệnh đề nào sau đây là đúng?
	 \choice
		{$B\Rightarrow A$} 
		{$B\Leftrightarrow  A$}
		 {$\overline{B}\Leftrightarrow\overline{A}$}
		{\True  $B\Rightarrow\overline{A}$}
	\loigiai{ 
		Ta có $ A $ sai nên $\overline{A}$ đúng, suy ra $B\Rightarrow\overline{A}$ đúng. 
	 }
\end{ex}

% nguồn: MaHoaNH_TracNghiem_Lop10
\begin{ex}%[0D1K1-4]%[ma:aeaee5b6]
	Cho hai mệnh đề $P$ và $Q$. Mệnh đề $P\Leftrightarrow Q$ đúng  khi
	\choice 
		{$P$  đúng và $Q$ sai}
		{$\overline{P}$ đúng  và $Q$ đúng}
		{$P$ sai  và $P$ đúng}
		{\True $\overline{P}$  đúng và $\overline{Q}$ đúng}
	 \loigiai{
		Ta có $\overline{P}$ đúng và $\overline{Q}$ đúng, suy ra $P$ sai và $P$  sai nên $P\Leftrightarrow Q$ đúng.
	} 
\end{ex}

% nguồn: MaHoaNH_TracNghiem_Lop10
\begin{ex}%[0D1B1-4]%[ma:bc8dab95]
	Trong các mệnh đề sau,  mệnh đề nào sai?
	\choice 
		{\True Điều kiện đủ để số tự nhiên $n$ chia  hết cho $24$ là $n$ chia hết cho $6$ và $4$}
		{Điều kiện đủ để $n^2+20$ là một hợp số  là $n$ là một số nguyên tố lớn hơn $3$}
		{Điều kiện đủ để $n^2-1$ chia hết cho $24$ là  $n$ là một số nguyên tố lớn hơn $3$}
		{Điều kiện đủ để một số nguyên dương tận cùng bằng $5$ là  số đó chia hết cho $5$}
	 \loigiai{
		Với $n=12$ thì  $n$ chia hết cho $6$ và $4$ nhưng $n$ không chia hết cho $24$
	} 
\end{ex}

% nguồn: MaHoaNH_TracNghiem_Lop10
\begin{ex}%[0D1B1-4]%[ma:6ce86626]
	Trong các mệnh đề  sau, mệnh đề nào đúng?
	\choice 
		{Nếu  $a\geq b$ thì $a^2\geq b^2$}
		{Nếu $a^2\geq b^2$ thì $a\geq b$ } 
		{\True  Nếu $a$ chia hết cho $9$ thì $a$ chia hết cho $3$}
		{Nếu $a$ chia hết cho $3$ thì $a$ chia hết cho $9$} 
	 \loigiai{
		Ta có $ a\,\vdots\,9 $ suy  ra $ a\,\vdots\,3 $.
	} 
\end{ex}

% nguồn: MaHoaNH_TracNghiem_Lop10
\begin{ex}%[0D1K1-4]%[ma:47ca5556]
	Cho các mệnh đề  $P: $ “$\left(\sqrt{125}+\sqrt{5}\right)^2$ là số nguyên” và
	$Q: $ “$\exists x \in  \mathbb{Q}: x-2=0$”. Phát biểu nào sau đây đúng?
	 \choice
		{$P \Rightarrow Q$ là mệnh  đề sai}
		{\True $\overline{Q} \Rightarrow  P$ là mệnh đề đúng}
		 {$\overline{P} \Rightarrow \overline{Q}$ là mệnh đề sai}
		{$P \Leftrightarrow \overline{Q}$ là mệnh đề  đúng}
	\loigiai{Ta có $\left(\sqrt{125}+\sqrt{5}\right)^2=180 \in \mathbb{Z}$ nên $P$ đúng. Dễ thấy  $Q$ đúng suy ra $\overline{Q}$ sai.\\
	Từ  đó $\overline{Q} \Rightarrow P$ là mệnh đề đúng.}
\end{ex}

% nguồn: MaHoaNH_TracNghiem_Lop10
\begin{ex}%[0D1B1-4]%[ma:a438884e]
	Cho $P$ là mệnh đề đúng, $Q$ là mệnh đề sai. Mệnh đề nào sau đây {\bf sai}? 
	\choice 
		 {$\overline{Q}$}
		 {$Q \Rightarrow P$}
		{\True  $P \Leftrightarrow Q$}
		{$P \Leftrightarrow  \overline{Q}$}
	\loigiai{ 
		Mệnh đề $Q$ và  $\overline{Q}$ là hai mệnh đề phủ định của nhau. $Q$ là mệnh đề sai thì $\overline{Q}$ là mệnh đề đúng.\\
	Mệnh đề $P \Leftrightarrow Q$ chỉ sai khi $P$ đúng và $Q$ sai.} 
\end{ex}

% nguồn: MaHoaNH_TracNghiem_Lop10
\begin{ex}%[0D1K1-4]%[ma:081df916]
	Trong các  mệnh đề sau, mệnh đề nào đúng?
	 \choice
		{Điều kiện  cần và đủ để tứ giác là hình thoi là khi có thể nội tiếp trong tứ giác đó một đường tròn}
		{\True Với các số thực dương $a$ và $b$, điều kiện cần và đủ để  $\sqrt{a}+\sqrt{b}=\sqrt{2(a+b)}$ là $a=b$}
		{Điều kiện cần và đủ để hai số tự nhiên dương $m$ và $n$ đều không chia  hết cho $9$ là $mn$ không chia hết cho $9$ }
		{Điều kiện cần và đủ để hai tam giác bằng nhau là hai tam giác đồng  dạng}
	\loigiai{ 
		Mệnh đề:  \lq\lq Với các số thực dương $a$ và $b$, điều kiện cần và đủ để $\sqrt{a}+\sqrt{b}=\sqrt{2(a+b)}$ là $a=b$ \rq\rq\, là mệnh đề đúng. Thật vậy:\\
		Với mọi số thực dương $a$ và $b$ giả sử $a=b$ thì $\heva{&\sqrt{a}+\sqrt{b}=\sqrt{a}+\sqrt{a}=2\sqrt{a}\\ &\sqrt{2(a+b)}=\sqrt{2(a+a)}=2\sqrt{a}}$ từ đó suy ra $\sqrt{a}+\sqrt{b}=\sqrt{2(a+b)}$. \\Ngược lại, giả sử $\sqrt{a}+\sqrt{b}=\sqrt{2(a+b)}\Rightarrow a+2\sqrt{a}\sqrt{b}+b=2(a+b)\Rightarrow  \left(\sqrt{a}-\sqrt{b}\right)^2=0\Rightarrow \sqrt{a}=\sqrt{b}$ hay $a=b$.
	 }
\end{ex}

% nguồn: MaHoaNH_TracNghiem_Lop10
\begin{ex}%[0D1K1-4]%[ma:09a9441f]
	Cho $P \Leftrightarrow Q$  là mệnh đề đúng. Khẳng định nào sau đây là \textbf{sai}?
	 \choice
		{$\overline{P} \Leftrightarrow Q$ sai} 
		{$\overline{P} \Leftrightarrow \overline{Q}$  đúng}
		{$\overline{Q} \Leftrightarrow P$ sai} 
		 {\True $\overline{P} \Leftrightarrow \overline{Q}$ sai}
	 \loigiai{
		Ta có $P \Leftrightarrow Q$ đúng nên $P \Rightarrow Q$ đúng và $Q \Rightarrow P$ đúng.\\ 
		Do đó $\overline{P} \Rightarrow \overline{Q}$ đúng và  $\overline{Q} \Rightarrow \overline{P}$ đúng.\\
	Vậy $\overline{P} \Leftrightarrow \overline{Q}$  đúng.}
\end{ex}

% nguồn: MaHoaNH_TracNghiem_Lop10
\begin{ex}%[0D1B1-4]%[ma:cdf9e8e8]
	Trong các mệnh đề sau đây, mệnh đề nào có mệnh đề  đảo đúng?
	\choice 
		{Nếu cả hai số chia hết cho $3$ thì tổng của  hai số đó chia hết cho $3$}
		{Nếu hai tam giác bằng nhau thì chúng  có diện tích bằng nhau}
		{Nếu một  số tự nhiên có chữ số tận cùng bằng $0$ thì số đó chia hết cho $5$ }
		{\True Nếu một số chia hết cho $5$ thì nó có chữ  số tận cùng bằng $0$}
	\loigiai{ 
		Xét mệnh đề $P\Rightarrow Q\colon$  ``Nếu một số chia hết cho $5$ thì nó có chữ số tận cùng bằng $0$''.\\
		Ta có mệnh đề đảo $ Q\Rightarrow P\colon $ ``Nếu một số  có chữ số tận cùng bằng $0$ thì số đó chia hết cho 5'' là mệnh đề đúng.
	 }
\end{ex}

% nguồn: MaHoaNH_TracNghiem_Lop10
\begin{ex}%[0D1B1-4]%[ma:c76887d6]
	Cho định lí \lq\lq  Nếu hai tam giác bằng nhau thì diện tích chúng bằng nhau\rq\rq. Mệnh đề nào sau đây đúng?
	 \choice
		{\True Hai tam  giác bằng nhau là điều kiện cần để diện tích chúng bằng nhau}
		{Hai tam giác bằng nhau là điều  kiện cần và đủ để chúng có diện tích bằng nhau}
		{Hai tam giác có diện tích bằng nhau  là điều kiện đủ để chúng bằng nhau}
		{Hai tam giác bằng nhau là điều kiện đủ  để diện tích chúng bằng nhau}
	\loigiai{ Với mệnh đề kéo theo \lq\lq Nếu hai tam giác bằng nhau thì diện  tích chúng bằng nhau\rq\rq. Thì ta nói\\
		\lq\lq Hai tam giác bằng nhau \rq\rq\, là điều kiện đủ.\\ 
		\lq\lq  Diện tích bằng nhau \rq\rq\, là điều kiện cần.
	 }
\end{ex}

% nguồn: MaHoaNH_TracNghiem_Lop10
\begin{ex}%[0D1K1-4]%[ma:3000c51e]
	Trong các mệnh  đề sau, mệnh đề nào \textbf{không phải} là định lí?
	\choice 
		{$\forall x \in \mathbb{N}, x^2 $ chia hết cho $3 \Rightarrow x$ chia hết cho $3$} 
		{$\forall x \in \mathbb{N}, x^2 $ chia  hết cho $6 \Rightarrow x$ chia hết cho $3$}
		{\True $\forall x \in \mathbb{N}, x^2 $ chia hết cho $ 9\Rightarrow x$ chia hết cho $9$} 
		{$ \forall x \in \mathbb{Z}, x$ chia hết cho $4$ và $6 \Rightarrow x$ chia hết  cho $12$}
	\loigiai{ 
		Xét  mệnh đề $\forall x \in \mathbb{N}, x^2 $ chia hết cho $9 \Rightarrow x$ chia hết cho $9$, với $x=3$ thì $x^2=3^2=9$ chia hết cho $9$, nhưng $3$ không chia hết cho $9$.\\
		Do đó mệnh đề $\forall x \in \mathbb{N}, x^2 $ chia  hết cho $ 9\Rightarrow x$ chia hết cho $9$ không phải là định lí.
	 }
\end{ex}

% nguồn: MaHoaNH_TracNghiem_Lop10
\begin{ex}%[0D1B1-4]%[ma:e4d6b301]
	Cho $P$ là  mệnh đề đúng, $Q$ là mệnh đề sai. Chọn mệnh đề đúng trong các mệnh đề sau.
	 \choice
		 {$P \Rightarrow \overline{P}$}
		 {$P \Leftrightarrow Q$}
		{\True  $\overline{P \Rightarrow Q}$}
		{$\overline{Q} \Rightarrow \overline{P}$} 
	 \loigiai{
		Ta có  $P$ là mệnh đề đúng, $Q$ là mệnh đề sai nên mệnh đề $P \Rightarrow Q$ là mệnh đề sai.\\
		Do đó $\overline{P \Rightarrow Q}$ là mệnh  đề đúng.
	} 
\end{ex}

% nguồn: MaHoaNH_TracNghiem_Lop10
\begin{ex}%[0D1K1-4]%[ma:5660c339]
	Cho $A$ là mệnh đề đúng, $B$ là mệnh đề sai, $C$ là  mệnh đề đúng. Mệnh đề nào sau đây {\bf sai}?
	 \choice
		{$A\Rightarrow C$} 
		 {$C\Rightarrow \left(A\Rightarrow\overline{B}\right)$}
		{$\left(\overline{B}\Rightarrow C\right)\Rightarrow  A$}
		{\True $C\Rightarrow (A\Rightarrow B)$} 
	 \loigiai{
		Ta có mệnh đề  $P\Rightarrow Q$ chỉ sai khi $P$ đúng $Q$ sai.\\
		 Từ giả thiết ta có $C$ đúng và $A\Rightarrow B$ sai. Vậy $C\Rightarrow (A\Rightarrow B)$ sai.
	} 
\end{ex}

% nguồn: MaHoaNH_TracNghiem_Lop10
\begin{ex}%[0D1B1-4]%[ma:0bb64b6e]
	 Trong các mệnh đề sau, mệnh đề nào \textbf{sai}?
	\choice 
		{\True $  -\pi<-2\Leftrightarrow \pi ^2<4$}
		{$ \pi<4\Leftrightarrow\pi ^2<16$} 
		{$\sqrt{23}<5\Rightarrow 2\sqrt{23}<2\cdot 5  $}
		{$\sqrt{23}<5\Rightarrow (-2)\sqrt{23}>(-2)\cdot 5  $}
	 \loigiai{
		Ta có $\pi ^2<4\Leftrightarrow  |\pi|<2\Leftrightarrow-2<\pi<2$.\\
		Vậy phương  án $ -\pi<-2\Leftrightarrow \pi ^2<4$ sai.
	} 
\end{ex}

% nguồn: MaHoaNH_TracNghiem_Lop10
\begin{ex}%[0D1K1-4]%[ma:6387c17a]
	Cho các mệnh đề $P,Q,R$ trong đó $R$ là mệnh đề đúng. Gọi $x,y$ là  giá trị của các mệnh đề $P,Q$,
	$x,y$ nhận các giá trị đúng hoặc sai. Có  tất cả bao nhiêu cặp giá trị $(x;y)$ sao cho
	mệnh đề $(R \Rightarrow  P) \Leftrightarrow (R \Rightarrow Q)$ đúng?
	 \choice
		{1} 
		{\True  2}
		{3} 
		{4} 
	\loigiai{Mệnh đề $(R \Rightarrow P) \Leftrightarrow (R \Rightarrow Q)$ đúng hai mệnh đề $R \Rightarrow P$ và $R \Rightarrow$  cùng đúng hoặc cùng sai.\\
		Ta có $R$ là mệnh đề đúng, nên $R \Rightarrow P$ và $R \Rightarrow$ cùng đúng khi hai mệnh đề $P$ và  $Q$ cùng đúng.\\
	Hai mệnh đề $R \Rightarrow P$ và $R \Rightarrow$ cùng sai khi hai  mệnh đề $P$ và $Q$ cùng sai.}
\end{ex}

% nguồn: MaHoaNH_TracNghiem_Lop10
\begin{ex}%[0D1B1-4]%[ma:d6776eeb]
	Trong  các mệnh đề sau, mệnh đề nào đúng?
	\choice 
		{ \True Điều kiện cần và đủ để mỗi số nguyên $a, b$ chia hết cho $7$ là tổng các bình phương của  chúng chia hết cho $7$}
		{Để tổng hai số tự nhiên chia hết cho $7$, điều kiện cần và  đủ là mỗi số đó chia hết cho $7$ }
		{Điều kiện cần và đủ để hai số nguyên dương  $a$ và $b$ đều không chia hết cho $9$ là tích $a\cdot b$ không chia hết cho $9$}
		{Để $a\cdot b >0$, điều kiện cần và đủ là hai số $a$ và $b$  đều dương}
	\loigiai{ 
		Với $ a +b\ \vdots\ 7$ thì chưa kết luận được $a\ \vdots\ 7$ và $b\ \vdots\ 7$ nên đáp án  B sai.\\
		Với $ a=3, b= 6$ không chia hết cho $9$ nhưng $a\cdot b =  18\ \vdots\ 9$ nên đáp án C sai.\\
		 Với $ a<0, b<0$ nhưng $a\cdot b >0$ nên đáp án D sai.\\
	 }
\end{ex}

% nguồn: MaHoaNH_TracNghiem_Lop10
\begin{ex}%[0D1B1-4]%[ma:b93ce247]
	Trong các mệnh đề sau, mệnh đề nào có  mệnh đề đảo đúng?
	\choice 
		{Nếu $a<b$ thì  $a\cdot t<b\cdot t$}
		{Nếu $x>y$  thì $x^2>y^2$}
		{Nếu $x=y$ thì  $t\cdot x=t\cdot y$}
		{\True  Nếu $a=b$ thì $a+c=b+c$}
	 \loigiai{
		Ta có 
		\begin{itemize}
			\item Mệnh đề \lq\lq Nếu $a\cdot t<b\cdot  t$ thì $a<b$\rq\rq\, chỉ đúng khi các số $a,b$ không âm và $t>0$.
			\item Mệnh đề \lq\lq Nếu $x^2>y^2$  thì $x>y$\rq\rq\, chỉ đúng khi $x,y\geq 0$.
			\item Nếu $t=0$ thì mệnh đề \lq\lq Nếu $t\cdot x=t\cdot y$ thì $x=y$\rq\rq\,  là không đúng.
			\item Mệnh đề  \lq\lq Nếu $a+c=b+c$ thì $a=b$\rq\rq\, đúng với mọi $a,\,b,\,c$.
	\end{itemize}}
\end{ex}

% nguồn: MaHoaNH_TracNghiem_Lop10
\begin{ex}%[0D1K1-4]%[ma:26e6fab1]
	Trong các mệnh đề  sau, mệnh đề nào là định lí?
	\choice 
		 {$\forall x \in \mathbb{R}, x>-2 \Rightarrow x^2>4$}
		{\True $\forall x  \in \mathbb{R}, x>2 \Rightarrow x^2>4$}
		 {$\forall x \in \mathbb{R}, x^2>4 \Rightarrow x>2$}
		{Nếu $a+b$ chia  hết cho $3$ thì $a,b$ đều chia hết cho $3$}
	 \loigiai{
		\begin{itemize}
			\item Xét mệnh đề $\forall x  \in \mathbb{R}, x>-2 \Rightarrow x^2>4$: Với $x=1>-2$ nhưng $(-1)^2=1<4$. Do đó mệnh đề này sai.
			\item Xét mệnh đề $\forall x \in \mathbb{R}, x^2>4 \Rightarrow x>2$, ta có $x^2>4\Rightarrow |x|>2\Rightarrow x<-\sqrt{2}$ hoặc $x>\sqrt{2}$. Do đó mệnh đề này sai. 
			\item Xét mệnh đề \lq\lq Nếu $a+b$ chia hết cho $3$ thì $a,b$ đều chia hết  cho $3$\rq\rq, ta chọn $a=5, b=1$ thì $a+b=6$ chia hết cho $3$ nhưng $a$ và $b$ đều không chia hết cho $3$. Do đó mệnh đề này sai.
		\end{itemize}
		Vậy mệnh  đề $\forall x \in \mathbb{R}, x>2 \Rightarrow x^2>4$ là định lí.
	 }
\end{ex}

% nguồn: MaHoaNH_TracNghiem_Lop10
\begin{ex}%[0D1B1-4]%[ma:71df2be9]
	 Trong các phát biểu mệnh đề $P\Rightarrow Q$ cho dưới đây, phát biểu nào {\bf sai}?
	 \choice
		{Nếu $P$ thì  $Q$}
		{$P$ kéo theo  $Q$}
		{\True $P$ là điều kiện cần  để có $Q$}
		{$P$  là điều kiện đủ để có $Q$}
	\loigiai{Mệnh đề $P\Rightarrow Q$ có thể phát biểu lại là: 
		\begin{itemize}
			\item Nếu  $P$ thì $Q$.
			\item $P$ kéo  theo $Q$.
			\item $P$ là điều kiện đủ để có $Q$. 
			\item $Q$ là điều kiện cần để có  $P$.
		\end{itemize}
	} 
\end{ex}

% nguồn: MaHoaNH_TracNghiem_Lop10
\begin{ex}%[0D1B1-4]%[ma:79b43f5e]
	Trong các  mệnh đề sau, mệnh đề nào là đúng?
	 \choice
		 {Tổng của hai số tự nhiên là một số chẵn khi và chỉ khi cả hai số đều là số chẵn}
		 {Tích của hai số tự nhiên là một số chẵn khi và chỉ khi cả hai số đều là số chẵn}
		{Tổng của hai số tự nhiên là một số  lẻ khi và chỉ khi cả hai số đều là số lẻ}
		{\True Tích của hai số tự nhiên là một số  lẻ khi và chỉ khi cả hai số đều là số lẻ}
	 \loigiai{
		Ta có các nhận  xét sau đây:
		\begin{itemize}
			\item Hai số lẻ vẫn  có tổng là số chẵn.
			\item Chỉ cần ít  nhất một thừa số chẵn thì tích hai số là số chẵn.
			\item Tổng hai số tự nhiên lẻ lại là một số  chẵn.
			\item Tích hai số  tự nhiên lẻ thì trong các thừa số không có thừa số chẵn, hay nói cách khác tất cả các thừa số đều là số lẻ.
		\end{itemize}
	 }
\end{ex}

% nguồn: MaHoaNH_TracNghiem_Lop10
\begin{ex}%[0D1K1-4]%[ma:bc05cb44]
	Cho các mệnh đề  $P$ đúng, $Q$ đúng, $R$ sai. Có bao nhiêu mệnh đề đúng trong các
	mệnh  đề sau ?\\
	1) $(P \Rightarrow Q) \Rightarrow  R$.\\
	2) $\overline{R} \Rightarrow (P\Rightarrow  Q)$.\\
	3) $(P \Rightarrow R) \Leftrightarrow Q$.\\ 
	4) $(\overline{R} \Leftrightarrow Q) \Leftrightarrow P$. 
	 \choice
		{1} 
		{\True  2}
		{3} 
		{4} 
	\loigiai{Ta có mệnh đề $P \Rightarrow Q$ đúng và $R$ sai nên mệnh đề $(P \Rightarrow Q)  \Rightarrow R$ sai.\\
		Ta có mệnh đề  $\overline{R}$ đúng và mệnh đề $P \Rightarrow Q$ đúng nên mệnh đề $\overline{R} \Rightarrow (P\Rightarrow Q)$ đúng.\\
		Mệnh đề $P \Rightarrow R$  là sai nên mệnh đề $(P \Rightarrow R) \Leftrightarrow Q$ là sai.\\
	Mệnh đề  $\overline{R}$ là sai nên $\overline{R} \Leftrightarrow Q$ là đúng dó đó mệnh đề $(\overline{R} \Leftrightarrow Q) \Leftrightarrow P$ đúng.}
\end{ex}

% nguồn: MaHoaNH_TracNghiem_Lop10
\begin{ex}%[0D1B1-4]%[ma:8dc79db0]
	Cho $a$, $b$  là hai số thực. Xét hai mệnh đề
	\begin{enumerate}[(1)]
		\item Nếu $a>b$ thì $a^2>b^2$. 
		 \item Nếu $a^2>b^2$ thì $a>b$.
	\end{enumerate}
	Khẳng  định nào sau đây là khẳng định đúng?
	\choice 
		{Chỉ $(1)$ đúng} 
		 {Chỉ $(2)$ đúng}
		{Cả $(1)$ và $(2)$  đều đúng}
		{\True Cả $(1)$ và $(2)$  đều sai}
	 \loigiai{
		Ta  có cả $2$ mệnh đề trên đều sai.\\
		Thật vậy mệnh đề $(1)$  sai vì với $a=1>b=-2$ nhưng $a^2=1<b^2=4$.\\
		Mệnh đề $(2)$ sai vì với $(-2)^2=4  >(-1)^2=1$ nhưng $-2< -1$.
	} 
\end{ex}

% nguồn: MaHoaNH_TracNghiem_Lop10
\begin{ex}%[0D1B1-4]%[ma:510fbf66]
	Trong các mệnh  đề sau, mệnh đề nào đúng?
	 \choice
		{\True $\forall x\in \mathbb{R}, x\le x^2  $}
		{ $\forall  x\in \mathbb{R}, |x|<3\Leftrightarrow x<3 $}
		{$\forall n\in  \mathbb{N}, n^2+1 $ chia hết cho $3$}
		{$\exists a\in \mathbb{Q}, a^2=2$} 
	 \loigiai{
		\begin{itemize}
			\item Mệnh đề $\forall x\in \mathbb{R},  |x|<3\Leftrightarrow x<3 $ sai vì $-4<3$ nhưng $|-4|>3$.
			 \item Mệnh đề $\forall n\in \mathbb{N}, n^2+1 $ chia hết cho $3$ sai, vì chẳng hạn chọn $n=1\in \mathbb{N}$ thì $2$ không chia hết cho $3$.
			 \item Xét mệnh đề $\exists a\in \mathbb{Q}, a^2=2$. Ta có
			\[a^2=2 \Leftrightarrow a=\pm \sqrt{2}  \in \mathbb{I}.\]
			Do  đó mệnh đề này sai.
		\end{itemize}
		Vậy mệnh  đề đúng là $\forall x\in \mathbb{R}, x\le x^2 $.
	} 
\end{ex}

% nguồn: MaHoaNH_TracNghiem_Lop10
\begin{ex}%[0D1K1-4]%[ma:8c6365e5]
	Giải bài toán \lq\lq chứng minh rằng, với các số $x,\,y,\,z$ bất  kỳ thì các bất đẳng thức sau không đồng thời xảy ra $\left|x\right|<\left|y-z\right|$; $\left|y\right|<\left|z-x\right|$; $\left|z\right|<\left|x-y\right|$\rq\rq. Một học sinh lập luận tuần tự như sau:
	\begin{enumerate}[Bước 1.]
		\item Giả sử các bất đẳng  thức trên xãy ra đồng thời.
		\item Thế thì nâng lên bình phương  hai vế các bất đẳng thức, chuyển vế phải sang trái, rồi phân tích, ta được:\\
		$(x-y+z)(x+y-z)<0$; $(y-z+x)(y+z-x)<0$; $(z-x+y)(z+x-y)<0$. 
		\item Sau đó nhân vế theo vế của các bất đẳng  thức trên thì thu được\\
		 $(x-y+z)^2(x+y-z)^2(-x+y+z)^2<0$, vô lý.
		\item Vậy các bất đẳng thức trên không thể  đồng thời xảy ra.
	\end{enumerate}
	Khẳng  định nào sau là đúng?
	 \choice
		 {Lí luận trên sai từ bước $2$}
		{Lí  luận trên sai từ bước $3$}
		{Lí luận trên sai từ bước  $2$}
		{\True Lí luận trên là lí luận  đúng}
	\loigiai{Lí luận trên là một lí luận hợp lí và hoàn toàn chính xác cho bài toán này. Các bước giải đều đúng với tính chất  của bất đẳng thức dẫn tới điều vô lí ở bước $3$. Tức là nếu các hệ thức trên đồng thời xảy ra thì dẫn tới điều vô lí trong toán học.
	} 
\end{ex}

% nguồn: MaHoaNH_TracNghiem_Lop10
\begin{ex}%[0D1B1-4]%[ma:b0182ed8]
	Cho các  mệnh đề sau:
	\begin{multicols}{2}
		\begin{enumerate}
			 \item ``Tam giác $ABC$ đều'';
			\item ``Tam giác $ABC$  có 2 góc bằng $60^\circ$'';
			\item ``Tam giác $ABC$  cân'';
			 \item ``$\triangle ABC$ có ba cạnh $AB=AC=BC$''.
		\end{enumerate}
	\end{multicols}
	Hỏi có bao nhiêu cặp mệnh đề  tương đương?
	\choice 
		{\True $3$} 
		{$2$} 
		{$6$} 
		 {$4$}
	\loigiai{ 
		Ta có $ P\Leftrightarrow Q $, $ Q\Leftrightarrow S $ và  $ S\Leftrightarrow P $.
	} 
\end{ex}

% nguồn: MaHoaNH_TracNghiem_Lop10
\begin{ex}%[0D1K1-4]%[ma:14f43bd8]
	Giải bài toán \lq\lq Cho $m$ là một số nguyên. Chứng minh  nếu $n^2$ chia hết cho $3$ thì $n$ chia hết cho $3$\rq\rq. Một học sinh làm như sau:
	\begin{enumerate}[Bước 1.]
		\item Giả sử $m$ không chia hết cho  $3$. Thế thì $m$ có một trong hai dạng $m=3k+1$ hoặc $m=3k+1$, với $k\in\mathbb{Z}$.
		\item Nếu $m=3k+1$  thì $m^2=9k^2+6k+1=3(3k^2+2k)+1$, còn nếu $m=3k+2$ thì $m^2=9k^2+12k+4=3(3k^2+4k+1)+1$.
		\item Vậy trong cả hai  trường hợp $m^2$ cũng không chia hết cho $3$, trái với giả thiết.
		\item Do đó $m$ là số lẻ nên phải chia hết  cho $3$.
	\end{enumerate}
	Lý luận trên đúng  tới bước nào?
	\choice 
		{\True  Bước $1$}
		{Bước $2$} 
		{Bước $3$} 
		{\True Bước  $4$}
	\loigiai{Lí luận trên sai từ bước $4$, vì tới đó không khẳng định được $m$ là số lẻ mà $m$ chỉ chia cho  $3$ dư $1$. Tức là nó không chia hết cho $3$.\\
		 Các bước giải đầy đủ là
		\begin{enumerate}[Bước 1.]
			\item Giả sử $m$ không chia hết cho $3$. Thế thì $m$ có một trong hai dạng $m=3k+1$ hoặc $m=3k+1$, với  $k\in\mathbb{Z}$.
			\item Nếu $m=3k+1$  thì $m^2=9k^2+6k+1=3(3k^2+2k)+1$, còn nếu $m=3k+2$ thì $m^2=9k^2+12k+4=3(3k^2+4k+1)+1$.
			\item Vậy trong cả hai trường hợp $m^2$  cũng không chia hết cho $3$, trái với giả thiết.
			 \item Do đó $m$ phải chia hết cho $3$.
		\end{enumerate}
	} 
\end{ex}

% nguồn: MaHoaNH_TracNghiem_Lop10
\begin{ex}%[0D1K1-4]%[ma:5603c9ba]
	Giải bài toán \lq\lq chứng minh với mọi số tự nhiên $n$ nếu $n^2$ là số chẵn  thì $n$ là số chẵn\rq\rq. Một bạn làm như sau:
	\begin{enumerate}[Bước 1.]
		\item Giả sử $n^2$  là chẵn.
		\item Thế thì $n$ có  dạng $n=2k+1$ với $k\in\mathbb{N}$.
		\item Khi đó  $n^2=(2k+1)^2=4k^2+4k+1=2(2k^2+2k)+1$.
		\item Ta thấy $2(2k^2+2k)+1$ là số lẻ nên $n^2$ là số lẻ. Trái với giả  thiết.
	\end{enumerate}
	Lý luận trên sai từ bước  nào?
	\choice 
		{\True Bước $1$} 
		 {Bước $2$}
		 {Bước $3$}
		{Bước $4$} 
	\loigiai{Mệnh đề có dạng  $P\Rightarrow Q$ nên giả thiết đã cho $n^2$ là số chẵn, ta cần giả sử $n$ là số lẻ để suy ra mâu thuẫn.
		Các bước giải  đầy đủ là
		\begin{enumerate}[Bước 1.]
			\item Giả sử $n$ là  số lẻ.
			\item Thế thì $n$ có dạng $n=2k+1$ với $k\in\mathbb{N}$. 
			\item Khi đó $n^2=(2k+1)^2=4k^2+4k+1=2(2k^2+2k)+1$. 
			\item Ta thấy $2(2k^2+2k)+1$ là số lẻ nên $n^2$ là số lẻ. Trái với  giả thiết.
		\end{enumerate}
	 }
\end{ex}


%%%%%%%%%%%%%%%% ĐÚNG/SAI %%%%%%%%%%%%%%%%

% nguồn: MaHoaNH_DungSai_Lop10
\begin{ex}%[0D1N1-4]%[ma:967874b0]
	Cho hai mệnh đề $P$: \lq\lq Tứ giác $ABCD$ là hình vuông\rq\rq\ và mệnh đề $Q$: \lq\lq Tứ giác $ABCD$ là hình chữ nhật có hai đường chéo  vuông góc với nhau\rq\rq. Các câu sau là đúng hay sai?
	 \choiceTF
		{\True Mệnh đề đảo của mệnh đề \lq\lq $P \Rightarrow Q$ là mệnh đề \lq\lq Nếu $ABCD$ là hình chữ nhật có hai đường chéo vuông góc  với nhau thì tứ giác $ABCD$ là hình vuông\rq\rq}
		{Hai mệnh đề $P$ và $Q$ không tương đương  với nhau}
		{Mệnh đề $P \Leftrightarrow Q$ là  mệnh đề sai}
		{\True $P$ là điều  kiện cần và đủ để có $Q$}
	\loigiai{ 
	\begin{itemchoice}
		\itemch Mệnh đề đảo của mệnh đề \lq\lq $P \Rightarrow Q$ là mệnh đề \lq\lq Nếu $ABCD$ là hình chữ nhật có hai đường chéo vuông góc với nhau thì tứ giác  $ABCD$ là hình vuông\rq\rq.
		\itemch Hai mệnh  đề $P$ và $Q$ tương đương với nhau.
		 \itemch Mệnh đề $P \Leftrightarrow Q$ là mệnh đề đúng.
		\itemch Vì $P$ và $Q$ tương đương nên $P$ là điều kiện cần và đủ  để có $Q$.
	\end{itemchoice}
	} 
\end{ex}

% nguồn: MaHoaNH_DungSai_Lop10
\begin{ex}%[0D1N1-4]%[ma:9d40a135]
	Cho hai mệnh đề $P$: \lq\lq$ABC$ là tam giác đều\rq\rq\ và mệnh đề $Q$: \lq\lq$ABC$ là tam giác cân\rq\rq. Xét mệnh đề kéo  theo \lq\lq Nếu $P$ thì $Q$\rq\rq. Các câu sau là đúng hay sai?
	 \choiceTF
		{\True $P$ là điều kiện đủ  để có $Q$}
		 {\True $Q$ là điều kiện cần để có $P$}
		 {\True Mệnh đề \lq\lq Nếu $P$ thì $Q$\rq\rq\ là mệnh đề đúng}
		{Mệnh  đề \lq\lq Nếu $Q$ thì $P$\rq\rq\ là mệnh đề đúng}
	\loigiai{ 
	\begin{itemchoice}
		\itemch $P$ là điều kiện đủ để có  $Q$.
		\itemch $Q$ là  điều kiện cần để có $P$.
		\itemch Mệnh đề \lq\lq Nếu $P$ thì $Q$\rq\rq\ là mệnh  đề đúng.
		\itemch Mệnh đề \lq\lq Nếu $Q$ thì  $P$\rq\rq\ là mệnh đề sai.
	\end{itemchoice}
	} 
\end{ex}


%%%%%%%%%%%%%%%% TỰ LUẬN %%%%%%%%%%%%%%%%

% nguồn: MaHoaNH_TuLuan_Lop10
\begin{bt}%[0D1H1-4]%[ma:3fa23dcd]
	Lập mệnh đề $P \Rightarrow Q$ và xét tính đúng  sai của nó, với $P$: \lq\lq $\pi > 4$\rq\rq\, và $Q$: \lq\lq$\pi^2 > 10$\rq\rq.
	\loigiai{ 
		Ta có mệnh đề $P \Rightarrow Q$ là \lq\lq Nếu $\pi >  4$ thì $\pi^2 > 10$\rq\rq.\\
		Vì $P$ sai (và $Q$ sai) nên mệnh đề $P  \Rightarrow Q$ là mệnh đề đúng.
	} 
\end{bt}

% nguồn: MaHoaNH_TuLuan_Lop10
\begin{bt}%[0D1H1-4]%[ma:a6449d1a]
	Phát biểu mệnh đề đảo của mệnh đề \lq\lq Nếu $\widehat{A}=90^\circ$ thì $\triangle ABC$ là tam giác vuông\rq\rq\, và  xét tính đúng sai của mệnh đề đó.
	 \loigiai{
		Ta có mệnh đề $P \Rightarrow Q$: \lq\lq Nếu $\widehat{A}=90^\circ$  thì $\triangle ABC$ là tam giác vuông\rq\rq.\\
		Mệnh đề đảo của mệnh đề trên là $Q \Rightarrow P$: \lq\lq Nếu $\triangle ABC$ là tam giác  vuông thì $\widehat{A}=90^\circ$\rq\rq.\\
		Mệnh đề $Q \Rightarrow P$ là mệnh đề sai, ví dụ trường hợp $\triangle ABC$ vuông  tại $B$.
	 }
\end{bt}

% nguồn: MaHoaNH_TuLuan_Lop10
\begin{bt}%[0D1H1-4]%[ma:b6f50320]
	Phát biểu mệnh đề $P \Leftrightarrow Q$  và xét tính đúng sai của nó. Giải thích
	\begin{itemize}
		 \item $P$: \lq\lq Bất phương trình $x^2 - 3x + 1 > 0$ có nghiệm\rq\rq.
		\item $Q$: \lq\lq Bất phương trình $x^2 - 3x + 1 \leq 0$ vô nghiệm\rq\rq. 
	\end{itemize}
	 \loigiai{
		Mệnh đề $P \Leftrightarrow Q$: \lq\lq Bất phương trình $x^2 - 3x + 1 > 0$ có  nghiệm khi và chỉ khi bất phương trình $x^2 - 3x + 1 \leq 0$ vô nghiệm\lq\lq.\\
		Mệnh đề trên sai. Vì bất phương  trình $x^2 - 3x + 1 \leq 0$ có nghiệm $x=1$.
	 }
\end{bt}

% nguồn: MaHoaNH_TuLuan_Lop10
\begin{bt}%[0D1H1-4]%[ma:f28e66b1]
	Cho hai mệnh đề $P$  và $Q$:
	\begin{itemize}
		\item $P$: \lq\lq $ABCD$ là tứ giác  nội tiếp\rq\rq.
		\item $Q$:  \lq\lq Tổng số đo hai góc đối nhau bằng $180^\circ$\rq\rq.
	\end{itemize}
	Hãy phát biểu  mệnh đề $P \Rightarrow Q$ dưới dạng điều kiện cần và điều kiện đủ.
	 \loigiai{
		Điều kiện cần: \lq\lq $ABCD$ là tứ giác nội tiếp là điều kiện cần để tổng số đo hai góc  đối nhau bằng $180^\circ$\rq\rq.\\
		Điều  kiện đủ: \lq\lq Trong tứ giác $ABCD$, tổng số đo hai góc đối nhau bằng $180^\circ$ là điều kiện đủ để $ABCD$ là tứ giác nội tiếp\rq\rq.
	} 
\end{bt}

% nguồn: MaHoaNH_TuLuan_Lop10
\begin{bt}%[0D1H1-4]%[ma:8544e19a]
	Phát biểu mệnh đề đảo của mệnh  đề sau và xét tính đúng sai của mệnh đề đảo.
	\begin{enumerate}
		\item Nếu một số chia hết cho $6$ thì  số đó chia hết cho $3$.
		\item Nếu một số là số tự nhiên  lẻ thì nó là số nguyên tố.
	\end{enumerate}
	\loigiai{ 
		\begin{enumerate}
			\item Nếu một số chia  hết cho $3$ thì số đó chia hết cho $6$. Đây là mệnh đề sai.\\
			Vì $ 9 $ chia hết cho $  3 $ nhưng $ 9 $ không chia hết cho $ 6 $.
			\item Nếu một số là số nguyên tố thì nó là số lẻ. Đây là mệnh đề sai vì $2$ là  số nguyên tố chẵn.
		\end{enumerate}
	} 
\end{bt}

% nguồn: MaHoaNH_TuLuan_Lop10
\begin{bt}%[0D1H1-4]%[ma:94ee4a68]
	Phát biểu mệnh đề $P \Leftrightarrow Q$ bằng hai cách và xét tính đúng  sai của nó.
	\begin{listEX}
		\item $P$: \lq\lq Tứ giác $ABCD$ là hình thoi\rq\rq\, và $Q$: \lq\lq Tứ giác $ABCD$ là hình bình hành  có hai đường chéo vuông góc với nhau\rq\rq.
		\item $P$: \lq\lq Bất phương trình $\sqrt{x^2 - 3x} > 1$ có nghiệm\rq\rq\, và  $Q$: \lq\lq$\sqrt{(-1)^2 - 3\cdot(-1)} > 1$\rq\rq.
	\end{listEX}
	\loigiai{ 
		\begin{listEX}
			\item Ta có mệnh đề $P \Leftrightarrow Q$ đúng vì mệnh đề $P \Rightarrow  Q$, $Q \Rightarrow P$ đều đúng và được phát biểu bằng hai cách như sau:\\
			\lq\lq Tứ giác $ABCD$ là  hình thoi khi và chỉ khi tứ giác $ABCD$ là hình bình hành có hai đường chéo vuông góc với nhau\rq\rq\, và
			\lq\lq Tứ giác $ABCD$ là hình thoi nếu và chỉ nếu tứ giác $ABCD$ là hình bình hành có hai đường chéo  vuông góc với nhau\rq\rq.\\
			\item Ta có mệnh đề $P \Leftrightarrow Q$ đúng vì mệnh đề $P, Q$  đều đúng (do mệnh đề $P \Rightarrow Q$, $Q \Rightarrow P$ đều đúng) và được phát biểu bằng hai cách như sau:\\
			\lq\lq Bất phương trình $\sqrt{x^2 - 3x} > 1$ có nghiệm khi và chỉ khi $\sqrt{(-1)^2 - 3\cdot(-1)} > 1$\rq\rq\, và 
			\lq\lq Bất phương trình $\sqrt{x^2 - 3x} > 1$ có nghiệm  nếu và chỉ nếu $\sqrt{(-1)^2 - 3\cdot(-1)} > 1$\rq\rq.
		\end{listEX}
	 }
\end{bt}

% nguồn: MaHoaNH_TuLuan_Lop10
\begin{bt}%[0D1H1-4]%[ma:df6d0f85]
	Hãy xác định giả thiết, kết luận đồng thời dùng thuật ngữ \lq\lq điều  kiện đủ\rq\rq\, để phát biểu các định lí sau
	\begin{enumerate}
		\item Nếu $a$ và $b$ là hai số hữu  tỉ thì tổng $a+b$ cũng là số hữu tỉ.
		\item Nếu  một số tự nhiên $n$ có tổng các chữ số chia hết cho $9$ thì nó chia hết cho $9$.
	\end{enumerate}
	 \loigiai{
		\begin{enumerate}
			\item Giả thiết của định lí là \lq\lq $a$ và $b$ là hai số  hữu tỉ\rq\rq.\\
			Kết  luận của định lí là \lq\lq tổng $a+b$ là số hữu tỉ\rq\rq.\\
			Phát biểu định lí dưới dạng điều kiện đủ \lq\lq Điều kiện đủ để tổng $a+b$ là số hữu tỉ là cả hai số $a$ và $b$ đều  là số hữu tỉ\rq\rq.
			\item Giả thiết của định lí là  \lq\lq Một số tự nhiên $n$ có tổng các chữ số chia hết cho $9$\rq\rq.\\
			Kết luận của định lí là \lq\lq $n$ chia hết cho $9$\rq\rq.\\ 
			Phát biểu định lí dưới dạng điều kiện đủ \lq\lq Điều kiện đủ để $n$ chia hết cho $9$ là tổng các chữ số của  $n$ chia hết cho $9$\rq\rq.
		\end{enumerate}
	 }
\end{bt}

% nguồn: MaHoaNH_TuLuan_Lop10
\begin{bt}%[0D1H1-4]%[ma:5c08c907]
	Cho $n$ là số tự nhiên. Xét  các mệnh đề:
	\begin{itemize}
		\item  $P\colon$ \lq\lq$n$ là một số tự nhiên chia hết cho $16$\rq\rq.
		\item $Q\colon$ \lq\lq$n$  là một số tự nhiên chia hết cho $8$\rq\rq.
	\end{itemize}
	Hãy phát biểu các mệnh đề sau \lq\lq $P\Rightarrow Q$\rq\rq, \lq\lq  $Q\Rightarrow P$\rq\rq, \lq\lq $P\Leftrightarrow Q$\rq\rq \, và xét tính đúng sai của các mệnh đề đó.
	 \loigiai{
		\begin{itemize}
			\item Phát biểu mệnh đề $P\Rightarrow Q\colon$\lq\lq Nếu $n$ là một số  tự nhiên chia hết cho $16$ thì $n$ là một số tự nhiên chia hết cho $8$\rq\rq.\\
			Mệnh đề này đúng, vì $n$ chia  hết cho 16 thì $n=16\cdot k\left(k\in \mathbb{N}\right)$ thì $n=8\cdot (2k)$ chia hết cho $8$.
			\item Phát biểu  mệnh đề $Q\Rightarrow P\colon$\lq\lq Nếu $n$ là một số tự nhiên chia hết cho $8$ thì $n$ là một số tự nhiên chia hết cho $16$\rq\rq \\
			Mệnh đề  này sai, chẳng hạn $n=8$ là số tự nhiên chia hết cho 8 nhưng $n$ không chia hết cho $16$.
			\item Do mệnh  đề \lq\lq$P\Rightarrow Q\colon$\rq\rq đúng, mệnh đề \lq\lq$Q\Rightarrow P\colon$\rq\rq sai nên mệnh đề \lq\lq $P\Leftrightarrow Q$\rq\rq sai.
		\end{itemize}
	} 
\end{bt}

% nguồn: MaHoaNH_TuLuan_Lop10
\begin{bt}%[0D1H1-4]%[ma:641a5945]
	 Phát biểu mệnh đề $P \Rightarrow Q$ và phát biểu mệnh đề đảo, xét tính đúng sai của nó.
	\begin{listEX}
		\item $P$: \lq\lq Tứ giác  $ABCD$ là hình thoi\rq\rq\, và $Q$: \lq\lq Tứ giác $ABCD$ có $AC$ và $BD$ cắt nhau tại trung điểm mỗi đường\rq\rq.
		\item $P$: \lq\lq$2 > 9$\rq\rq\, và $Q$: \lq\lq$4 <  3$\rq\rq.
		\item $P$: \lq\lq Tam  giác $ABC$ vuông cân tại $A$\rq\rq\, và $Q$: \lq\lq Tam giác $ABC$ có $\widehat{A}=2\widehat{B}$\rq\rq.
		\item $P$: \lq\lq Ngày 2 tháng 9 là ngày Quốc Khánh của nước Việt Nam\rq\rq\,  và $Q$: \lq\lq Ngày 27 tháng 7 là ngày thương binh liệt sĩ\rq\rq.
	\end{listEX}
	 \loigiai{
		\begin{listEX}
			\item Mệnh đề $P \Rightarrow Q$ là \lq\lq Nếu tứ giác $ABCD$ là hình thoi thì $AC$  và $BD$ cắt nhau tại trung điểm mỗi đường\rq\rq, mệnh đề này đúng.
			Mệnh đề đảo là $Q \Rightarrow P$: \lq\lq Nếu tứ giác $ABCD$ có $AC$  và $BD$ cắt nhau tại trung điểm mỗi đường thì $ABCD$ là hình thoi\rq\rq, mệnh đề này sai.
			\item  Mệnh đề $P \Rightarrow Q$ là \lq\lq Nếu $2 > 9$ thì $4 < 3$\rq\rq, mệnh đề này đúng vì mệnh đề $P$ sai.
			Mệnh đề đảo là  $Q \Rightarrow P$: \lq\lq Nếu $4 < 3$ thì $2 > 9$\rq\rq, mệnh đề này đúng vì mệnh đề $Q$ sai.
			\item Mệnh đề $P \Rightarrow Q$ là \lq\lq Nếu tam giác $ABC$ vuông cân tại $A$ thì $\widehat{A}=2\widehat{B}$\rq\rq, mệnh đề này  đúng.
			Mệnh đề đảo là $Q \Rightarrow P$: \lq\lq Nếu tam giác  $ABC$ có $\widehat{A}=2\widehat{B}$ thì nó vuông cân tại $A$\rq\rq, mệnh đề này sai.
			\item Mệnh đề $P \Rightarrow Q$ là \lq\lq Nếu ngày 2 tháng  9 là ngày Quốc Khánh của nước Việt Nam thì ngày 27 tháng 7 là ngày thương binh liệt sĩ\rq\rq.
			Mệnh đề đảo là $Q \Rightarrow P$: \lq\lq Nếu ngày 27 tháng 7 là ngày thương binh liệt sĩ thì ngày 2 tháng 9 là ngày Quốc Khánh của nước Việt Nam\rq\rq. Hai mệnh  đề trên đều đúng vì mệnh đề $P$, $Q$ đều đúng.
		\end{listEX}
	} 
\end{bt}

% nguồn: MaHoaNH_TuLuan_Lop10
\begin{bt}%[0D1B1-4]%[ma:8f7f3cd3]
	Sử dụng các thuật ngữ \lq\lq điều kiện cần\rq\rq, \lq\lq điều kiện đủ\rq\rq\, để phát biểu lại định lí: \lq\lq Nếu tứ giác $ABCD$ là hình chữ nhật thì hai đường  chéo bằng nhau\rq\rq.
	 \loigiai{
		Ta có thể phát biểu  lại định lí đã cho như sau:\\
	\lq\lq Tứ giác  $ABCD$ có hai đường chéo bằng nhau là điều kiện cần để nó là hình chữ nhật\rq\rq\, hoặc \lq\lq Tứ giác $ABCD$ là hình chữ nhật là điều kiện đủ để hai đường chéo bằng nhau\rq\rq.}
\end{bt}

% nguồn: MaHoaNH_TuLuan_Lop10
\begin{bt}%[0D1K1-4]%[ma:8bfb6725]
	Chứng minh bằng  phản chứng. Cho $n \in \mathbb{N}$, nếu $5n+5$ là số lẻ thì $n$ là số chẵn.
	 \loigiai{
		Giả sử $n$ là số lẻ. Suy ra $n=2k+1,~\forall  k \in \mathbb{N}$.\\
		Khi đó $5n+5=5(2k+1)+5=10k+10=10(k+1)$ (trái giả thiết $5n+5$ là  số lẻ).\\
		Vậy phải có  $n$ là số chẵn.
	 }
\end{bt}

% nguồn: MaHoaNH_TuLuan_Lop10
\begin{bt}%[0D1B1-4]%[ma:e060be7f]
	Lập mệnh đề kéo  theo và mệnh đề tương đương của hai mệnh đề sau đây và cho biết tính đúng, sai của chúng. Biết:\\
	$P$ : \lq\lq  Điểm $M$ nằm trên phân giác của góc $Oxy$\rq\rq.\\
	$Q$ : \lq\lq  Điểm $M$ cách đều hai cạnh $Ox$, $Oy$ \rq\rq.
	 \loigiai{
		\begin{itemize}
			\item Mệnh đề $P\Rightarrow Q$ là: \lq\lq Nếu một điểm bất kỳ nằm trên đường phân giác của góc $Oxy$  thì nó cách đều hai cạnh $Ox$ và $Oy$\rq\rq. Đây là mệnh đề đúng.\\
			\item Mệnh đề $P\Leftrightarrow Q$ là: \lq\lq Mọi điểm nằm trên đường phân giác của góc $Oxy$ khi và chỉ khi chúng cách đều hai cạnh $Ox$ và  $Oy$\rq\rq. Đây là mệnh đề đúng.
		\end{itemize}
	} 
\end{bt}

% nguồn: MaHoaNH_TuLuan_Lop10
\begin{bt}%[0D1B1-4]%[ma:a2b87dbf]
	Dùng thuật ngữ "điều kiện cần và đủ " đề phát biều  định lí sau
	\begin{enumerate}
		\item Một tam giác là tam giác cân nếu  và chỉ nếu nó có hai góc bằng nhau.
		\item Tứ giác là hình bình hành khi và chỉ khi tứ giác có hai đường chéo cắt nhau tại trung  điểm của mỗi đường.
	\end{enumerate}
	\loigiai{	\begin{enumerate}
		\item Điều kiện cần và đủ để một tam giác là  tam giác cân là nó có hai góc bằng nhau.
		\item Điều kiện cần  và đủ để rt giác là hình bình hành là tứ giác có hai đường chéo cắt nhau tại trung điểm của mỗi đường.
\end{enumerate}}
\end{bt}

% nguồn: MaHoaNH_TuLuan_Lop10
\begin{bt}%[0D1B1-4]%[ma:3f01d306]
	Dùng khái niệm " điều kiện cần "  để phát biểu các định lí sau
	\begin{enumerate}
		\item Nếu $M A \perp  M B$ thì $M$ thuộc đường tròn đường kính $A B$.
		 \item $a \neq 0$ hoặc $b \neq 0$ là điều kiện đủ để $a^{2}+b^{2}>0$.
	\end{enumerate}
	\loigiai{\begin{enumerate}
		\item Điều  kiện cần để $M A \perp M B$ là $M$ thuộc đường tròn đường kính $A B$.
		\item Điều  kiện cần để $a^{2}+b^{2}>0$ là $a \neq 0$ hoặc $b \neq 0$.
\end{enumerate}}
\end{bt}

% nguồn: MaHoaNH_TuLuan_Lop10
\begin{bt}%[0D1B1-4]%[ma:c4df6e9c]
	Trong mỗi cặp  mệnh đề $P$ và $Q$ sau đây, hãy phát biểu mệnh đề $P\Rightarrow Q$ và xét tính đúng sai của nó. $P$ có phải là điều kiện đủ để có $Q$ không?
	\begin{enumerate}
		\item $P\colon$ \lq\lq$a$ và $b$ là hai số chẵn\rq\rq \,  , $Q\colon$ \lq\lq $a+b$ là số chẵn\rq\rq \,($a$, $b$ là hai số tự nhiên).
		 \item $P\colon$ \lq\lq Tứ giác $ABCD$ có bốn cạnh bằng nhau\rq\rq \, , $Q\colon$ \lq\lq Tứ giác $ABCD$ là một hình vuông\rq\rq.
	\end{enumerate}
	\loigiai{\begin{enumerate}
		\item $P\Rightarrow Q$ : \lq\lq Nếu $a$ và $b$ là số chẵn thì $a+b$ là số chẵn\rq\rq.\\ Đây là mệnh đề {\bf đúng} nên $P$ là điều kiện đủ để có  $Q$.
		\item $P\Rightarrow Q$: \lq\lq Nếu tứ giác $ABCD$ có bốn cạnh bằng nhau thì tứ giác $ABCD$ là hình vuông\rq\rq.\\ Đây là mệnh đề {\bf sai} nên $P$ không là điều kiện đủ  để có $Q$.
\end{enumerate}}
\end{bt}

% nguồn: MaHoaNH_TuLuan_Lop10
\begin{bt}%[0D1B1-4]%[ma:08dd0738]
	Sử dụng thuật ngữ \lq\lq điều kiện  cần và đủ \rq\rq để phát biểu các định lí sau
	\begin{enumerate}
		\item Một tứ giác nội tiếp được trong một đường tròn khi và chỉ khi tổng hai góc đối diện của  nó bằng $180^{\circ}$.
		\item Tam giác cân khi và chỉ khi có trung tuyến bằng  nhau.
	\end{enumerate}
	\loigiai{	\begin{enumerate}
		\item Điều kiện cần và đủ  để một tứ giác nội tiếp được trong một đường tròn là tổng hai góc đối diện của nó bằng $180^{\circ}$.
		\item Điều kiện cần và đủ để tam  giác cân là có trung tuyến bằng nhau.
\end{enumerate}}
\end{bt}

% nguồn: MaHoaNH_TuLuan_Lop10
\begin{bt}%[0D1B1-4]%[ma:55e53a55]
	Xét hai mệnh  đề:\\
	$P\colon$  \lq\lq Tam giác $ABC$ vuông tại $A$\rq\rq;\\
	$Q\colon$ \lq\lq  Tam giác $ABC$ có $AB^2+AC^2=BC^2$\rq\rq.\\
	Hai mệnh đề $P$ và $Q$ có tương đương không? Nếu có, hãy phát biểu một định lí thể hiện điều này, trong đó có sử dụng thuật ngữ \lq\lq khi và chỉ khi\rq\rq\, hoặc \lq\lq  điều kiện cần và đủ\rq\rq.
	 \loigiai{
		Theo  định lí Pythagore, hai mệnh đề $P \Rightarrow Q$ và $Q \Rightarrow P$ đều đúng.\\
		Do  đó, $P$ và $Q$ là hai mệnh đề tương đương.\\
		Ta có thể phát biểu thành định lí  như sau:\\
		\lq\lq Tam giác $ABC$ vuông tại $A$ khi và chỉ khi $AB^2+AC^2=BC^2$\rq\rq\, hoặc \lq\lq Để tam giác $ABC$ vuông tại $A$, điều kiện cần và đủ là  $AB^2+AC^2=BC^2$\rq\rq.
	 }
\end{bt}

% nguồn: MaHoaNH_TuLuan_Lop10
\begin{bt}%[0D1B1-4]%[ma:379cd97b]
	Phát biểu một  \lq\lq điều kiện đủ\rq\rq
	\begin{enumerate}[a.]
		\item Để tứ  giác $ABCD$ là hình bình hành.
		\item Để tứ giác $ABCD$ là hình  chữ nhật.
	\end{enumerate}
	 \loigiai{
		\begin{enumerate}[a.]
			\item Để tứ giác $ABCD$ là hình bình hành, điều kiện đủ là tứ giác  $ABCD$ có hai đường chéo cắt nhau tại trung điểm của mỗi đường.
			\item Để tứ giác $ABCD$  là hình chữ nhật, điều kiện đủ là hình bình hành $ABCD$ có một góc vuông.
		\end{enumerate}
	} 
\end{bt}

% nguồn: MaHoaNH_TuLuan_Lop10
\begin{bt}%[0D1B1-4]%[ma:6f59ac58]
	Sử dụng khái niệm "điều kiện cần " đề  phát biều các định lí sau
	\begin{enumerate}
		\item Nếu một số tự nhiên chia hết cho 15 thì  nó chia hết cho 5.
		 \item Nếu $a=b$ thì $a^2=b^2$.
		\item Trong mặt phằng, nếu hai đường thằng phân biệt cùng  vuông góc với một đường thằng thứ ba thì hai đường thằng ấy song song với nhau.
	\end{enumerate}
	\loigiai{	\begin{enumerate}
		\item Điều kiện cần để một số tự nhiên  chia hết cho 15 là nó chia hết cho 5.
		\item  Điều kiện cần để $a=b$ là $a^2=b^2$.
		\item Trong mặt phằng, điều kiện cần để hai đường thằng phân biệt cùng vuông góc với một đường thằng thứ ba là hai  đường thằng ấy song song với nhau.
\end{enumerate}}
\end{bt}

% nguồn: MaHoaNH_TuLuan_Lop10
\begin{bt}%[0D1B1-4]%[ma:9d3b4287]
	Hãy phát biểu mệnh đề kéo theo $P \Rightarrow Q$, $\overline{Q} \Rightarrow P$ và xét  đúng sai của mệnh đề này.
	\begin{enumerate}
		\item Cho tứ giác $A B C D$ và hai mệnh đề $P$ : "Tổng hai góc đối cùa tứ  giác lồi bằng $180^\circ$" và $Q:$ " Tứ giác nội tiếp được đường tròn".
		\item $P: "\sqrt 2-\sqrt 3>-1 $" và  $Q: "(\sqrt 2-\sqrt 3)^2>(-1)^2$".
	\end{enumerate}
	\loigiai{	\begin{enumerate}
		\item $P \Rightarrow Q$: " Nếu tổng hai góc đối cùa tứ giác lồi bằng $180^\circ$ thì tứ giác đó nội tiếp  được đường tròn" - Đúng.\\
		$\overline{Q} \Rightarrow  P:$ "Nếu tứ giác không nội tiếp được đường tròn thì tổng hai góc đối của tứ giác lồi bằng $180^\circ$" - Sai.
		\item $P \Rightarrow Q$: "Nếu $\sqrt 2-\sqrt 3>-1 $ thì $(\sqrt 2-\sqrt 3)^2>(-1)^2$" - Sai.\\ 
		$\overline{Q} \Rightarrow P:$ "Nếu $(\sqrt 2-\sqrt 3)^2\leq (-1)^2$ thì $\sqrt 2-\sqrt 3>-1 $  - Đúng.
\end{enumerate}}
\end{bt}

% nguồn: MaHoaNH_TuLuan_Lop10
\begin{bt}%[0D1B1-4]%[ma:0d20dbd1]
	Sừ dụng khái niệm "điều kiện  đủ " đề phát biểu các định lí sau
	\begin{enumerate}
		\item Nếu $a$ và $b$ là hai số  hũu tỉ thì tổng $a+b$ là số hũu tỉ.
		\item Nếu hai tam giác bằng nhau thì chúng có diện  tích bằng nhau.
		 \item Nếu một số tự nhiên có chữ số tận cùng là chữ số 5 thì nó chia hết cho 5.
	\end{enumerate}
	\loigiai{\begin{enumerate}
		\item $a$  và $b$ là hai số hũu tỉ là điều kiện đủ để có tổng $a+b$ là số hũu tỉ.
		\item Hai tam giác bằng nhau là điều kiện đủ  để chúng có diện tích bằng nhau.
		\item Một số tự nhiên có chữ số tận cùng là chữ số 5 là điều kiện đủ để nó chia  hết cho 5.
\end{enumerate}}
\end{bt}

% nguồn: MaHoaNH_TuLuan_Lop10
\begin{bt}%[0D1B1-4]%[ma:c4367060]
	 Dùng thuật ngữ "điều kiện cần và đủ " đề phát biều định lí sau
	\begin{enumerate}
		\item Tam giác  $ABC$ vuông khi và chi khi $AB^2+AC^2=BC^2$.
		\item  Tứ giác là hình chũ nhật khi và chỉ khi nó có ba góc vuông.
		\item Tứ giác là nội tiếp  được trong đường tròn khi và chỉ khi nó có hai góc đối bù nhau.
		\item Một số chia hết cho 2 khi  và chỉ khi nó có chữ số tận cùng là số chẵn.
	\end{enumerate}
	\loigiai{	\begin{enumerate}
		\item Điều kiện cần và đủ để tam giác $ABC$ vuông  là $AB^2+AC^2=BC^2$.
		\item Điều kiện cần và  đủ để tứ giác là hình chũ nhật là nó có ba góc vuông.
		\item Điều kiện cần và đủ để tứ giác là nội tiếp được trong đường tròn là nó có hai góc  đối bù nhau.
		\item Điều kiện cần và đủ để  một số chia hết cho 2 là nó có chữ số tận cùng là số chẵn.
\end{enumerate}}
\end{bt}

% nguồn: MaHoaNH_TuLuan_Lop10
\begin{bt}%[0D1B1-4]%[ma:7f340218]
	Sử dụng thuật ngữ \lq\lq điều kiện cần\rq\rq, \lq\lq điều kiện đủ\rq\rq đề phát biều định  lí sau
	\begin{enumerate}
		\item Nếu một tứ giác là hình vuông thì nó có bốn cạnh bằng nhau. Có định lí  đảo của định lí trên không, vì sao?
		\item Nếu một tứ giác là hình thoi thì nó có hai đường chéo  vuông góc. Có định lí đảo của định lí trên không, vì sao?
	\end{enumerate}
	\loigiai{\begin{enumerate}
		\item Điều kiện cần để một tứ giác là hình vuông là nó  có bốn cạnh bằng nhau.\\
		Điều kiện đủ để tứ giác có  bốn cạnh bằng nhau là tứ giác đó là hình vuông.\\
		Không tồn tại định lí đảo  của định lí đã cho. Vì mệnh đề đảo của mệnh đề trên là một mệnh đề sai.
		\item Điều kiện cần để một tứ giác là hình thoi là nó có hai  đường chéo vuông góc.\\
		Điều kiện đủ để một tứ  giác có hai đường chéo vuông góc là tứ giác đó là hình thoi.\\
		Không tồn tại định lí đảo của  định lí đã cho. Vì mệnh đề đảo của mệnh đề trên là một mệnh đề sai.
\end{enumerate}}
\end{bt}

% nguồn: MaHoaNH_TuLuan_Lop10
\begin{bt}%[0D1B1-4]%[ma:0a1a369b]
	Cho số thực $x$. Xét các mệnh đề $P$ : \lq\lq $x^2=1$ \rq\rq và  $Q$ : \lq\lq $x=1$ \rq\rq
	\begin{enumerate}[a.]
		 \item Phát biểu mệnh đề $P\Rightarrow Q$ và mệnh đề đảo của nó.
		\item Xét tính đúng sai của hai mệnh đề trên. 
		\item Chỉ  ra một giá trị của $x$ để mệnh đề $P\Rightarrow Q$ sai.
	\end{enumerate}
	 \loigiai{
		\begin{enumerate}[a.]
			\item Phát biểu mệnh  đề $P\Rightarrow Q$: Nếu bình phương của một số bằng $1$ thì số đó bằng $1$.\\Phát biểu mệnh đề $Q\Rightarrow P$: Nếu một số bằng $1$ thì bình phương số đó bằng $1$.
			\item Mệnh đề $P\Rightarrow Q$ là mệnh đề sai.\\Mệnh đề $Q\Rightarrow P$  là mệnh đề đúng.
			 \item Với $x=-1$ thì mệnh đề $P\Rightarrow Q$ sai.
		\end{enumerate}
	} 
\end{bt}

% nguồn: MaHoaNH_TuLuan_Lop10
\begin{bt}%[0D1B1-4]%[ma:da18f623]
	Cho định lí \lq\lq Cho  số tự nhiên $n$, nếu $n^{5}$ chia hết cho 5 thì $n$ chia hết cho $5$\rq\rq. Đinh lí này được viết dưới dạng $P \Rightarrow Q$.
	\begin{enumerate}
		\item  Hãy xác định các mệnh đề $P$ và $Q$.
		\item Phát biểu định lí trên bằng cách dùng thuật ngữ \lq\lq điều kiện  cần\rq\rq.
		\item Phát biểu định lí trên bằng cách dùng thuật ngữ  \lq\lq điều kiện đủ\rq\rq.
		\item Hãy phát biểu định lí đảo (nếu có) của định lí trên rồi dùng các thuật ngữ "điều kiện cần và đủ" phát  biều gộp cả hai định lí thuận và đảo.
	\end{enumerate}
	\loigiai{\begin{enumerate}
		 \item $P:$ $n^{5}$ chia hết cho 5 và $Q:$ $n$ chia hết cho $5$.
		\item Cho  số tự nhiên $n$, điều kiện cần để có $n^{5}$ chia hết cho 5 là $n$ chia hết cho $5$
		\item Cho số tự nhiên $n$, $n^{5}$ chia hết cho 5  là $n$ chia hết cho $5$
		\item Định lí đảo: \lq\lq Cho số tự nhiên $n$, nếu $n$ chia hết cho  $5$ thì $n^5$ chia hết cho $5$.\\
		Cho số tự nhiên $n$, điều kiện cần và đủ để $n^{5}$ chia  hết cho 5 là $n$ chia hết cho $5$.
\end{enumerate}}
\end{bt}

% nguồn: MaHoaNH_TuLuan_Lop10
\begin{bt}%[0D1B1-4]%[ma:ff554130]
	Sử dụng khái niệm \lq\lq điều kiện cần\rq\rq\, hoặc \lq\lq  điều kiện đủ\rq\rq\, phát biểu các mệnh đề sau:
	\begin{enumerate}[a.]
		\item Hai tam giác bằng nhau có diện  tích bằng nhau.
		\item Số tự nhiên có chữ số tận  cùng là chữ số $5$ thì nó chia hết cho $5$.
		\item  Nếu $a=b$ thì $a^2=b^2$.
		\item Nếu $a+b>0$ thì trong hai số $a$ và $b$ lớn hơn  $0$.
	\end{enumerate}
	 \loigiai{
		\begin{enumerate}[a.]
			\item Hai tam giác bằng nhau là điều kiện đủ để hai tam giác đó có diện tích bằng nhau.\\Hai tam giác bằng nhau  là điều kiện cần để chúng bằng nhau.
			\item Số tự nhiên có chữ số tận cùng là chữ số $5$ là điều kiện đủ để nó chia hết  cho $5$.\\Số tự nhiên chia hết cho $5$ là điều kiện cần để nó có tận cùng là chữ số $5$
			\item Bình phương của hai số bằng nhau là điều kiện cần để hai số đó bằng nhau.\\Hai số bằng nhau  là điều kiện đủ để bình phương của chúng bằng nhau.
			\item Hai số $a$ và $b$ lớn hơn $0$ là điều kiện đủ để tổng của chúng lớn hơn $0$.\\Tổng của  hai số lớn hơn $0$ là điều kiện cần để hai số đó đều lớn hơn $0$.
		\end{enumerate}
	} 
\end{bt}

% nguồn: MaHoaNH_TuLuan_Lop10
\begin{bt}%[0D1K1-4]%[ma:d214f7bb]
	Chứng minh các mệnh đề  sau bằng phương pháp phản chứng
	\begin{enumEX}{1}
		\item Nếu $x\ne 1$ và  $y\ne 1$ thì $x+y-xy\ne 1$.
		\item Nếu bình phương của một số  tự nhiên $n$ là một số chẵn thì $n$ cũng là một số chẵn.
	\end{enumEX}
	 \loigiai{
		\begin{enumEX}{1}
			\item  Nếu $x\ne 1$ và $y\ne 1$ thì $x+y-xy\ne 1$.\\
			Giả sử $x+y-xy=1$, ta có $(y-1)(x-1)=0\Leftrightarrow x=1\text{hoặc}y=1$ (mẫu thuẫn  giả thiết $x\ne 1$ và $y\ne 1$).\\
			Suy ra điều giả sử  sai. Vậy ta có điều cần chứng minh.
			 \item Nếu bình phương của một số tự nhiên $n$ là một số chẵn thì $n$ cũng là một số chẵn.\\
			Giả sử $n$  là một số lẻ.\\
			Suy ra $n=2k+1$ ($k\in \mathbb{N}$), $n^2=4k^2+4k+1$ không chia hết cho $2$. Suy ra $n^2$ là số lẻ (mâu thuẫn  giả thiết $n^2$ là số chẵn).\\
			Suy ra điều  giả sử sai. Vậy ta có điều cần chứng minh.
		\end{enumEX}
	} 
\end{bt}

% nguồn: MaHoaNH_TuLuan_Lop10
\begin{bt}%[0D1B1-4]%[ma:25cc9001]
	Phát  biểu các mệnh đề sau với thuật ngũ \lq\lq điều kiện cần \rq\rq, \lq\lq điều kiện đủ\rq\rq
	\begin{enumerate}
		\item Nếu hai tam giác  bằng nhau thì chúng có diện tích bằng nhau.
		\item  Nếu số nguyên dương chia hết cho 6 thì chia hết cho 3.
		\item Nếu hình thang có hai đường chéo bằng nhau thì  nó là hình thang cân.
		\item Nếu tam giác $A B C$  vuông tai $A$ và $A H$ là đường cao thì $A B^{2}=B C \cdot B H$.
	\end{enumerate}
	\loigiai{	\begin{enumerate}
		\item Điều kiện  cần để hai tam giác bằng nhau là chúng có diện tích bằng nhau.\\
		Điều kiện đủ để hai  tam giác có diện tích bằng nhau là hai tam giác bằng nhau.
		\item Điều kiện cần để số  nguyên dương chia hết cho 6 là chia hết cho 3.\\
		Điều kiện đủ để  số nguyên dương chia hết cho 3 là nó chia hết cho 6.
		\item Điều kiện cần hình thang có hai đường chéo bằng nhau là nó là hình thang cân.\\ 
		Điều kiện đủ để hình thang  là hình thang cân là hình thang có hai đường chéo bằng nhau.
		\item Điều kiện cần để tam  giác $A B C$ vuông tai $A$ và $A H$ là đường cao là $A B^{2}=B C \cdot B H$.\\
		Điều kiện đủ để tam giác $ABC$ có $A B^{2}=B C \cdot B H$ là tam giác $A B C$ vuông  tai $A$ và $A H$ là đường cao.
\end{enumerate}}
\end{bt}

% nguồn: MaHoaNH_TuLuan_Lop10
\begin{bt}%[0D1B1-4]%[ma:caeae6e0]
	Xét  tính đúng sai của các mệnh đề sau:
	\begin{enumerate}[a)]
		\item $R\colon$ \lq\lq Nếu tam giác $ABC$ có hai  góc bằng $60^\circ$ thì nó là tam giác đều.\rq\rq;
		\item $T\colon$ \lq\lq Từ $-3<-2$ suy  ra $(-3)^2<(-2)^2$\rq\rq.
	\end{enumerate}
	\loigiai{ 
		\begin{enumerate}[a)]
			\item $R$ là mệnh  đề có dạng $P \Rightarrow Q$, với $P\colon$ \lq\lq Tam giác $ABC$ có hai góc bằng $60^\circ$\rq\rq\, và $Q\colon$ \lq\lq Tam giác $ABC$ là tam giác đều\rq\rq.
			Ta thấy khi $P$ đúng thi $Q$ cũng  đúng.\\
			Do  đó, $P \Rightarrow Q$ đúng.\\
			Vậy $R$ là  mệnh đề đúng.
			\item $T$ là mệnh đề có dạng $P \Rightarrow Q$, với $P\colon$ \lq\lq$-3<-2$\rq\rq\, và $Q\colon$  \lq\lq$(-3)^2<(-2)^2$\rq\rq\, (hay $Q\colon$ \lq\lq$9<4$\rq\rq).\\
			Ta thấy mệnh đề $P$ đúng, còn mệnh đề $Q$ sai.\\ 
			Do đó, $P  \Rightarrow Q$ sai.\\
			Vậy $T$ là  mệnh đề sai.
		\end{enumerate}
	 }
\end{bt}

% nguồn: MaHoaNH_TuLuan_Lop10
\begin{bt}%[0D1K1-4]%[ma:8e0902fb]
	Chứng minh các mệnh đề sau bằng phương  pháp phản chứng
	\begin{enumEX}{1}
		\item Nếu $a+b<2$ thì một trong hai số $a$ và $b$ nhỏ  hơn $1$.
		\item Một tam giác không phải là tam giác đều thì nó có ít nhất một góc  nhỏ hơn $60^\circ$.
	\end{enumEX}
	 \loigiai{
		\begin{enumEX}{1}
			\item Nếu $a+b<2$ thì một trong hai số $a$ và  $b$ nhỏ hơn $1$.\\
			Giả sử $a$ và $b$ đều lớn hơn  hoặc bằng $1$, suy ra $a+b\ge 2$ (mâu thuẫn giả thiết $a+b<2$).\\
			Suy ra điều giả  sử sai. Vậy nếu $a+b<2$ thì một trong hai số $a$ và $b$ nhỏ hơn $1$.
			\item Một tam giác không phải là tam giác đều thì nó có ít nhất một góc nhỏ  hơn $60^\circ$.\\
			Giả sử $\triangle ABC$ có  ba góc đều lớn hơn hoặc bằng $60^\circ$.\\
			Suy  ra $\widehat{A}+\widehat{B}+\widehat{C}\ge 180^\circ$. Mà $\widehat{A}+\widehat{B}+\widehat{C}= 180^\circ$ nên $\widehat{A}=\widehat{B}=\widehat{C}=60^\circ$.\\
			Suy ra $\triangle ABC$ đều (mâu thuẫn giả thiết $\triangle ABC$ không là tam giác  đều).\\
			Suy ra điều giả sử sai. Vậy một tam giác không phải là tam giác đều thì nó có ít  nhất một góc nhỏ hơn $60^\circ$.
		\end{enumEX}
	 }
\end{bt}

% nguồn: MaHoaNH_TuLuan_Lop10
\begin{bt}%[0D1B1-4]%[ma:a727ccb8]
	Phát biểu các mệnh đề sau bằng cách sử dụng khái niệm \lq\lq điều  kiện cần và đủ\rq\rq.
	\begin{enumEX}{1}
		\item Một tam  giác là vuông khi và chỉ khi nó có một góc bằng tổng hai góc còn lại.
		\item Một tứ giác là hình chữ nhật khi  và chỉ khi nó có ban góc vuông.
		\item Một tứ giác là  nội tiếp được trong đường tròn khi và chỉ khi nó có hai góc đối bù nhau.
		\item Một số chia hết cho $6$ khi và chỉ khi nó chia hết  cho $2$ và cho $3$.
		\item Số tự nhiên $n$  là số lẻ khi và chỉ khi $n^2$ là số lẻ.
	\end{enumEX}
	\loigiai{ 
		\begin{enumEX}{1}
			\item  Một tam giác là vuông khi và chỉ khi nó có một góc bằng tổng hai góc còn lại.\\
			Điều kiện cần và đủ để \lq\lq một tam giác là vuông\rq\rq\  là \lq\lq nó có một góc bằng tổng hai góc còn lại\rq\rq.
			\item Một tứ giác là hình chữ nhật khi và chỉ khi nó có ba góc vuông.\\ 
			Điều kiện cần và đủ để \lq\lq một tứ giác là  hình chữ nhật\rq\rq\ là \lq\lq nó có ba góc vuông\rq\rq.
			\item Một tứ giác là nội tiếp được trong đường tròn khi và chỉ khi nó có hai góc đối bù  nhau.\\
			Điều kiện cần và đủ để  \lq\lq một tứ giác là nội tiếp được trong đường tròn\rq\rq\ là \lq\lq nó có hai góc đối bù nhau\rq\rq.
			\item Một số  chia hết cho $6$ khi và chỉ khi nó chia hết cho $2$ và cho $3$.\\
			Điều kiện cần và đủ để \lq\lq một số chia hết cho $6$\rq\rq\ là \lq\lq nó chia hết cho $2$ và  cho $3$\rq\rq.
			\item Số  tự nhiên $n$ là số lẻ khi và chỉ khi $n^2$ là số lẻ.\\
			Điều kiện cần và đủ để \lq\lq  số tự nhiên $n$ là số lẻ\rq\rq\ là \lq\lq $n^2$ là số lẻ\rq\rq.
		\end{enumEX}
	 }
\end{bt}

% nguồn: MaHoaNH_TuLuan_Lop10
\begin{bt}%[0D1K1-4]%[ma:a5115a00]
	Chứng minh các mệnh đề  sau bằng phương pháp phản chứng
	\begin{enumEX}{1}
		\item Nếu tích của  hai số tự nhiên là một số lẻ thì tổng của chúng là một số chẵn.
		\item Nếu một tứ giác có tổng các góc đối diện bằng $2$ góc vuông thì  tứ giác nội tiếp được đường tròn.
		\item Nếu  $x^2+y^2=0$ thì $x=0$ và $y=0$.
	\end{enumEX}
	 \loigiai{
		\begin{enumEX}{1}
			\item  Nếu tích của hai số tự nhiên là một số lẻ thì tổng của chúng là một số chẵn.\\
			Giả sử tổng của hai số  tự nhiên $a$, $b$ là một số lẻ.\\
			Suy  ra $a$ chẵn, $b$ lẻ hoặc $a$ lẻ, $b$ chẵn.\\
			Suy ra  tích $ab$ là số chẵn (mâu thuẫn giả thiết tích $ab$ lẻ).\\
			Suy ra điều giả sử sai. Vậy ta  có điều cần chứng minh.
			\item Nếu một tứ giác có tổng các góc đối diện bằng $2$ góc vuông thì  tứ giác nội tiếp được đường tròn.\\
			Giả  sử tứ giác không nội tiếp.\\
			Suy ra tổng  các góc đối điện khác $180^\circ$ (khác $2$ góc vuông) (mâu thuẫn giả thiết).\\
			Suy ra điều giả sử sai. Vậy  ta có điều phải chứng minh.
			\item  Nếu $x^2+y^2=0$ thì $x=0$ và $y=0$.\\
			Giả  sử $x\ne 0$ hoặc $y\ne 0$.\\
			Suy ra $x^2>0$ hoặc  $y^2>0$.\\
			Suy ra $x^2+y^2>0$ (mẫu thuẫn giả thiết  $x^2+y^2=0$).\\
			Suy ra điều giả sử sai. Vậy ta có điều cần chứng  minh.
		\end{enumEX}
	} 
\end{bt}

% nguồn: MaHoaNH_TuLuan_Lop10
\begin{bt}%[0D1B1-4]%[ma:f68b4380]
	Các mệnh đề sau đúng hay sai? Nếu sai hãy sửa lại để có  một mệnh đề đúng.
	\begin{listEX}[2]
		\item  $x=1\Leftrightarrow x^2=1$.
		\item $2001$ là số  nguyên tố.
		\item $ABCD$ là hình vuông $\Rightarrow ABCD$ là hình  bình hành.
		\item $ABCD$  là hình thoi $\Rightarrow ABCD$ là hình chữ nhật.
		\item Tứ giác $MNPQ$ là hình vuông  $\Leftrightarrow$ Hai đường chéo $MP$ và $NQ$ bằng nhau.
		 \item Hai tam giác bằng nhau $\Leftrightarrow$ Chúng có diện tích bằng nhau.
	\end{listEX}
	 \loigiai{
		\begin{listEX}
			 \item $x=1\Leftrightarrow x^2=1$.\\
			Mệnh đề sai. Sửa lại thành  \lq\lq $x=1\Rightarrow x^2=1$\rq\rq.
			\item $2001$  là số nguyên tố.\\
			Mệnh đề sai. Sửa lại thành \lq\lq$2001$ không là số nguyên tố\rq\rq. 
			\item $ABCD$ là hình vuông $\Rightarrow ABCD$ là hình bình hành.\\ 
			 Mệnh đề đúng.
			\item $ABCD$ là hình thoi $\Rightarrow ABCD$ là hình chữ  nhật.\\
			Mệnh đề sai. Sửa lại thành \lq\lq $ABCD$ là hình thoi $\Rightarrow ABCD$ là hình bình  hành\rq\rq.
			\item Tứ giác $MNPQ$ là hình vuông $\Leftrightarrow$ Hai đường chéo $MP$ và  $NQ$ bằng nhau.\\
			Mệnh đề sai. Sửa lại thành \lq\lq Tứ giác $MNPQ$ là hình vuông $\Rightarrow$ Hai đường chéo $MP$ và  $NQ$ bằng nhau\rq\rq.
			\item Hai  tam giác bằng nhau $\Leftrightarrow$ Chúng có diện tích bằng nhau.\\
			Mệnh đề sai. Sửa lại thành \lq\lq Hai tam giác bằng nhau $\Rightarrow$ Chúng có diện  tích bằng nhau\rq\rq.
		\end{listEX}
	} 
\end{bt}

% nguồn: MaHoaNH_TuLuan_Lop10
\begin{bt}%[0D1B1-4]%[ma:4ed8f599]
	Phát biểu các mệnh đề sau bằng cách  sử dụng khái niệm \lq\lq điều kiện cần\rq\rq, \lq\lq điều kiện đủ\rq\rq.
	\begin{listEX}
		\item Trong mặt phẳng, nếu hai đường thẳng phân  biệt cùng vuông góc với một đường thẳng thứ ba thì hai đường thẳng ấy song song với nhau.
		\item Nếu hai tam giác bằng nhau thì chúng có diện tích bằng  nhau.
		\item Nếu tứ  giác $T$ là một hình thoi thì nó có hai đường chéo vuông góc với nhau.
		\item Nếu tứ giác $H$  là một hình chữ nhật thì nó có ba góc vuông.
		\item Nếu tam giác $K$ đều thì  nó có hai góc bằng nhau.
	\end{listEX}
	 \loigiai{
		\begin{listEX}
			\item Trong mặt phẳng, nếu hai đường thẳng phân biệt cùng vuông góc với một đường thẳng thứ  ba thì hai đường thẳng ấy song song với nhau.\\
			\lq\lq Trong mặt  phẳng, hai đường thẳng phân biệt song song với nhau\rq\rq\ là điều kiện cần để \lq\lq hai đường thẳng ấy cùng vuông góc với một đường thẳng thứ ba\rq\rq.\\
			\lq\lq Trong mặt phẳng, hai đường  thẳng phân biệt cùng vuông góc với một đường thẳng thứ ba\rq\rq\ là điều kiện đủ để \lq\lq hai đường thẳng ấy song song với nhau.
			\item  Nếu hai tam giác bằng nhau thì chúng có diện tích bằng nhau.\\
			\lq\lq  Hai tam giác có diện tích bằng nhau\rq\rq\ là điều kiện cần để \lq\lq chúng bằng nhau\rq\rq.\\
			\lq\lq Hai tam giác bằng nhau\rq\rq\ là  điều kiện đủ để \lq\lq chúng có diện tích bằng nhau\rq\rq.
			\item Nếu tứ giác $T$ là một hình thoi thì nó có hai đường chéo vuông góc  với nhau.\\
			\lq\lq Tứ  giác $T$ có hai đường chéo vuông góc với nhau\rq\rq\ là điều kiện cần để \lq\lq nó là một hình thoi\rq\rq.\\
			\lq\lq Tứ giác $T$ là một hình thoi\rq\rq\ là điều kiện đủ để \lq\lq nó có hai đường chéo vuông  góc với nhau\rq\rq.
			\item Nếu tứ giác $H$ là một hình chữ nhật thì nó có ba góc  vuông.\\
			\lq\lq Tứ giác $H$ có ba  góc vuông\rq\rq\ là điều kiện cần để \lq\lq nó là một hình chữ nhật\rq\rq.\\
			\lq\lq Tứ giác $H$ là một hình chữ nhật\rq\rq\ là điều kiện đủ để \lq\lq nó có ba góc  vuông\rq\rq.
			\item  Nếu tam giác $K$ đều thì nó có hai góc bằng nhau.\\
			\lq\lq Tam giác $K$ có hai góc bằng nhau\rq\rq\ là điều kiện cần  để \lq\lq nó đều\rq\rq.\\
			\lq\lq Tam giác $K$ đều\rq\rq\ là điều kiện đủ để \lq\lq nó  có hai góc bằng nhau\rq\rq.
		\end{listEX}
	} 
\end{bt}

% nguồn: MaHoaNH_TuLuan_Lop10
\begin{bt}%[0D1B1-4]%[ma:49735ffd]
	Phát biểu các mệnh đề sau, bằng cách sử dụng khái niệm \lq\lq điều kiện  cần\rq\rq, \lq\lq điều kiện đủ\rq\rq.
	\begin{enumEX}{1}
		\item Nếu một số tự nhiên có chữ số tận cùng là chữ số $5$ thì nó chia hết cho  $5$.
		\item Nếu $a+b>0$ thì một trong hai số $a$  và $b$ phải dương.
		\item Nếu một số tự nhiên chia hết cho $6$ thì  nó chia hết cho $3$.
		\item Nếu $a=b$  thì $a^2=b^2$.
		\item Nếu $a$ và $b$ cùng chia hết  cho $c$ thì $a+b$ chia hết cho $c$.
	\end{enumEX}
	 \loigiai{
		\begin{enumEX}{1}
			\item Nếu một số tự nhiên có chữ số tận cùng là chữ số $5$ thì nó chia hết cho $5$.\\ 
			\lq\lq Một số tự nhiên có chữ số tận cùng là chữ  số $5$\rq\rq\ là điều kiện đủ để \lq\lq nó chia hết cho $5$\rq\rq.\\
			 \lq\lq Một số tự nhiên chia hết cho $5$\rq\rq\ là điều kiện cần để \lq\lq nó có chữ số tận cùng là chữ số $5$\rq\rq.
			 \item Nếu $a+b>0$ thì một trong hai số $a$ và $b$ phải dương.\\
			\lq\lq $a+b>0$\rq\rq\ là điều kiện đủ để \lq\lq một  trong hai số $a$ và $b$ phải dương\rq\rq.\\
			\lq\lq Một trong hai số $a$ và $b$ phải dương\rq\rq\ là  điều kiện cần để \lq\lq $a+b>0$\rq\rq.
			\item Nếu một số tự nhiên chia hết cho $6$ thì nó chia  hết cho $3$.\\
			\lq\lq Một  số tự nhiên chia hết cho $6$\rq\rq\ là điều kiện đủ để \lq\lq nó chia hết cho $3$\rq\rq.\\
			\lq\lq Một số tự nhiên  chia hết cho $3$\rq\rq\ là điều kiện cần để \lq\lq nó chia hết cho $6$\rq\rq.
			\item Nếu $a=b$ thì  $a^2=b^2$.\\
			\lq\lq $a=b$\rq\rq\ là  điều kiện đủ để \lq\lq $a^2=b^2$\rq\rq.\\
			 \lq\lq $a^2=b^2$\rq\rq\ là điều kiện cần để \lq\lq $a=b$\rq\rq.
			\item Nếu $a$ và $b$ cùng chia hết cho $c$ thì $a+b$ chia hết cho $c$.\\ 
			\lq\lq $a$ và $b$ cùng chia hết cho $c$\rq\rq\ là điều kiện đủ để \lq\lq $a+b$ chia hết  cho $c$\rq\rq.\\
			\lq\lq $a+b$ chia hết cho $c$\rq\rq\ là điều kiện cần để \lq\lq $a$ và $b$ cùng chia hết  cho $c$\rq\rq.
		\end{enumEX}
	 }
\end{bt}
